\documentclass[11pt,a4paper]{article}
\usepackage[margin=2.5cm]{geometry}
\usepackage{amsmath,amssymb,amsthm}
\usepackage{enumitem}
\usepackage{mathrsfs}
\usepackage{url}

\newtheorem{theorem}{Theorem}[section]
\newtheorem{lemma}[theorem]{Lemma}
\newtheorem{proposition}[theorem]{Proposition}
\newtheorem{corollary}[theorem]{Corollary}
\newtheorem{remark}[theorem]{Remark}
\newtheorem{example}[theorem]{Example}

\newcommand{\h}{\mathfrak{h}}

\title{A uniform decomposition theorem for invariant differential operators on imprimitive complex reflection groups $G(r,p,n)$}
\author{Jean Kabor\'e$^1$, Ibrahim Nonkan\'e$^2$\\[4pt]
\small $^1$ Laboratoire de Sciences et Technologies (LaST),
Universit\'e Thomas Sankara, Burkina Faso\\
\small $^2$ D\'epart\'ement d'\'economie et de math\'ematiques appliqu\'ees, IUFIC,
Universit\'e Thomas Sankara, Burkina Faso\\[4pt]
\small \texttt{kaborejean775075@gmail.com}$^1$, \texttt{ibrahim.nonkane@uts.bf}$^2$}
\date{}

\begin{document}
\maketitle

\begin{abstract}
{We study the module structure of the polynomial ring,
localized at the discriminant, over the ring of differential
operators on the ring of invariants of the imprimitive complex
reflection group $G(r,p,n)$, describing its simple components with
explicit generators given by the higher Specht polynomials of
Ariki--Terasoma--Yamada. The proof rests on a Jacobian lemma computing
the discriminant of $G(r,p,n)$, combined with a double-centralizer
argument. As particular cases ($r=2$, $p=2$ or $p=1$) we recover, and
considerably shorten, the known decomposition theorems for the real
reflection groups $W(D_n)$ and $W(B_n)$; we also treat $G(r,r,n)$ and
$G(r,1,n)$ explicitly, with worked examples ($D_2$, $D_3$, $B_2$) and
their central idempotents. Finally, applying the Galois descent equivalence of categories of
Nonkan\'e to $G(r,p,n)$ for the first time
gives a second, generator-free description of
the simple summands as twisted invariants.}
\end{abstract}

\noindent\textbf{2020 Mathematics Subject Classification.}
Primary 20F55, 13A50; Secondary 16S32, 20C15, 05E10.

\noindent\textbf{Keywords.} Complex reflection groups, imprimitive
reflection groups, invariant theory, Weyl algebra, ring of invariant
differential operators, higher Specht polynomials, Young tableaux,
decomposition theorem, Galois descent, Coxeter groups of type $B$
and $D$.

\section{Introduction}\label{sec:intro}

The interplay between the invariant theory of a finite reflection
group $W$ acting on a polynomial ring
$\mathcal{O}_X = \mathbb{C}[x_1,\dots,x_n]$ and the representation
theory of $W$ itself is one of the oldest and most productive themes
in algebra, going back to the classical Chevalley--Shephard--Todd
theorem and its companion fact on coinvariant algebras: the invariant
subring $\mathcal{O}_X^W$ is a polynomial ring, and $\mathcal{O}_X$,
viewed as a module over it, is free of rank $|W|$, with fiber
isomorphic to the regular representation of $W$. When $W$ is a
\emph{complex} reflection group -- in particular one of the
imprimitive groups $G(r,p,n)$ in the Shephard--Todd classification --
a finer combinatorial model for this regular-representation structure
is furnished by the \emph{higher Specht polynomials} of Ariki,
Terasoma, and Yamada~\cite{ATY1997}, later extended to the full
family $G(r,p,n)$ by Morita and Yamada~\cite{MoritaYamada1998}.

A parallel and largely independent line of work studies
$\mathcal{O}_X$ not merely as a $\mathbb{C}[W]$-module, but as a
module over $\widetilde{\mathcal{D}}_Y$, the ring of differential
operators on the invariant subring $\mathcal{O}_Y = \mathcal{O}_X^W$
after localizing at the discriminant $\Delta$; this ring plays the
role of an invariant differential operator ring for $W$. Write
$\widetilde{\mathcal{O}}_X$ for the corresponding localization of
$\mathcal{O}_X$ at $\Delta$. Understanding the resulting
$\widetilde{\mathcal{D}}_Y$-module structure of $\widetilde{\mathcal{O}}_X$
-- its simple constituents, their multiplicities, and explicit
generators -- was carried out for the symmetric group $W = \mathcal{S}_n$
in~\cite{Nonkane2019}, extended to certain products of symmetric
groups in~\cite{NonkaneTodjihounde2023}, to the wreath product
$G(r,n) = \mathbb{Z}/r\mathbb{Z} \wr \mathcal{S}_n$ (that is, $G(r,1,n)$
in the notation below) in~\cite{NonkaneLawson2022}, and to the real
reflection group $W(D_n)$ in~\cite{NonkaneLawson2021} -- our own
starting point for the present generalization -- in all cases via
an equivalence of categories of Galois-descent type. The present paper
completes this picture by treating the \emph{entire} family of
imprimitive complex reflection groups $G(r,p,n)$ at once, rather than
case by case.

\paragraph{Main results.} Our first main result
(Theorem~\ref{thm:main-general}, Proposition~\ref{prop:decomposition-general})
identifies, for every $G(r,p,n)$, the simple
$\widetilde{\mathcal{D}}_Y$-submodules of $\widetilde{\mathcal{O}}_X$
and their multiplicities, with explicit generators given by the
higher Specht polynomials. The proof combines two ingredients that are
new to this generality: a single \emph{Jacobian lemma}
(Lemma~\ref{lem:jacobian-general}) computing the discriminant of
$G(r,p,n)$ as the determinant of the change of variables between
$\widetilde{\mathcal{O}}_X$ and $\widetilde{\mathcal{O}}_Y$, and a
self-contained \emph{double-centralizer argument}
(Lemma~\ref{lem:double-centralizer}) replacing the ad hoc, group-specific
constructions used previously. Because the argument is uniform in
$(r,p,n)$, specializing it to the two families of \emph{real}
reflection groups in this classification -- $W(D_n) = G(2,2,n)$ and
$W(B_n) = G(2,1,n)$ -- yields their decomposition theorems as
two-line corollaries (Corollaries~\ref{cor:Dn-main}
and~\ref{cor:Bn}), considerably shortening what a direct,
group-by-group treatment would require. We also work out the
remaining structural cases $G(r,r,n)$ and $G(r,1,n)$ explicitly, and
illustrate all of this on the smallest nontrivial examples
($D_2$, $D_3$, $B_2=G(2,1,2)$), including their central and primitive
idempotents.\footnote{Every explicit symbolic computation in these
examples -- the Jacobian lemma on a sample of $(n,r,p)$, the
generators and characters for $D_2$ and $D_3$, the central idempotents
for $D_2$, $D_3$, and $G(r,1,2)$, and the reflection counts for
$G(r,r,n)$ and $W(B_n)$ -- is reproduced and checked independently in
an accompanying SymPy script, publicly available at
\url{PLACEHOLDER-REPOSITORY-URL} (archived with a permanent DOI at
\url{PLACEHOLDER-DOI}).}

Our second main result (Subsection~\ref{subsec:galois-descent}) applies
the Galois-descent equivalence of categories of
Nonkan\'e~\cite[\S 2.4]{Nonkane2019}
(originally applied to the symmetric group in~\cite{Nonkane2019}, and
recalled and applied to products of symmetric groups
in~\cite{NonkaneTodjihounde2023}) to $W = G(r,p,n)$ for
the first time, and derives from it two new consequences: a second,
generator-free construction of the simple summands of
$\widetilde{\mathcal{O}}_X$ as ``twisted invariants''
(Corollary~\ref{cor:twisted-invariants}), and a categorical criterion
for recognizing when an abstract $\widetilde{\mathcal{D}}_Y$-module
occurs as a direct summand of $\widetilde{\mathcal{O}}_X$
(Corollary~\ref{cor:characterization}), without needing an explicit
higher Specht polynomial generator.

\paragraph{Structure of the paper.} Section~\ref{sec:prelim} fixes
notation for reflections and discriminants, and recalls the
combinatorics of higher Specht polynomials for $G(r,p,n)$. Section~\ref{sec:general}
contains our main results: the Jacobian lemma, the double-centralizer
argument, the decomposition theorem, and the Galois-descent
equivalence of categories. Section~\ref{sec:Dn} specializes to the
real reflection group $W(D_n)$, with the examples $D_2$ and $D_3$
worked out in full detail. Section~\ref{sec:other-cases} treats the
remaining cases $G(r,r,n)$, $W(B_n) = G(2,1,n)$, and $G(r,1,n)$,
again with explicit low-rank examples.

\section{Preliminaries}\label{sec:prelim}

\subsection{Reflections and the discriminant}\label{subsec:notation}

We first fix notation used throughout. Let $\h$ be a finite-dimensional
vector space. We say that a semisimple element $s \in GL(\h)$ is a
\emph{reflection} (or \emph{complex reflection}, when we wish to
stress the contrast with the real reflections defined next) if
$\mathrm{rank}(1-s)=1$. For the purposes of the next definition only,
assume in addition that
$\h$ carries a non-degenerate inner product $(\cdot,\cdot)$ and is
defined over $\mathbb{R}$. We say that a semisimple element
$s \in GL(\h)$ is a \emph{real reflection} if $\mathrm{rank}(1-s)=1$
and $s$ is an involution, i.e.\ $s^2 = \mathrm{id}$.\footnote{\label{fn:real}This
distinguishes a general (complex) reflection, which may have any
finite order, from a real reflection, of order exactly $2$, acting on
a real inner product space -- consistent with the terminology
``real reflection group'' used below for $W(D_n)$, as opposed to the
complex reflection groups $G(r,p,n)$ of Subsection~\ref{subsec:specht},
whose reflections have order $r$ in general.}

Let $s \in GL(\h)$ be a complex reflection. Denote by
$\alpha_s \in \h^*$ an eigenvector of $s$ in $\h^*$ with nontrivial
eigenvalue. For $W$ a reflection group with set of reflections
$S \subset W$, the \emph{polynomial discriminant} of $W$ is
$\Delta = \prod_{s \in S} \alpha_s(x)$ (well defined up to a nonzero
scalar, since each $\alpha_s$ is). This is the object whose
localization underlies the whole of Section~\ref{sec:general}: for a
general complex reflection group $G(r,p,n)$, $\Delta$ is computed
explicitly in Lemma~\ref{lem:jacobian-general} below, and specialized
to $W(D_n) = G(2,2,n)$ in Corollary~\ref{cor:Dn-setup}.

\subsection{Higher Specht Polynomials for the reflection group $G(r,p,n)$}\label{subsec:specht}

For a general complex reflection group $G(r,p,n)$ -- not assumed real
-- the $\widetilde{\mathcal{D}}_Y$-module structure of
$\widetilde{\mathcal{O}}_X$ studied in this paper is described using
the combinatorial machinery of \emph{higher Specht polynomials}, which
we now recall (see~\cite{MoritaYamada1998}); it underlies the whole of
Section~\ref{sec:general}. The group $G(r,p,n)$, where
$r,p,n \ge 1$ and $p$ divides $r$, consists of the monomial matrices
whose nonzero entries are of the form $\xi^j$ ($0 \le j < r$) and such
that the $d$-th power of the product of all nonzero entries is equal
to $1$, where $\xi$ denotes a primitive $r$-th root of unity and
$d = r/p$. The imprimitive complex reflection group $G(r,p,n)$ is
identified as a normal subgroup of the wreath product
\[
  G(r,n) = (\mathbb{Z}/r\mathbb{Z})^n \rtimes \mathcal{S}_n
    = \{ (\xi^{i_1},\dots,\xi^{i_n};\sigma) \mid i_k \in \mathbb{N},\ \sigma \in \mathcal{S}_n \}
  \;=\; G(r,1,n)
\]
(the case $p=1$, where the defining condition above is vacuous), where
$\mathcal{S}_n$ is the group of permutations of
$\{x_1,\dots,x_n\}$, whose product is given by
\[
  (\xi^{i_1},\dots,\xi^{i_n};\sigma)(\xi^{j_1},\dots,\xi^{j_n};\pi)
  = (\xi^{i_1+j_{\sigma^{-1}(1)}},\dots,\xi^{i_n+j_{\sigma^{-1}(n)}};\sigma\pi).
\]

Let $\mathcal{O}_X = \mathbb{C}[x_1,\dots,x_n]$ be the ring of
polynomials in $n$ indeterminates, on which $G(r,n)$ acts by
\[
  (\xi^{i_1},\dots,\xi^{i_n};\sigma) f
  = f(\xi^{i_{\sigma(1)}}x_{\sigma(1)},\dots,\xi^{i_{\sigma(n)}}x_{\sigma(n)}).
\]
It is well known that the fundamental invariants under this action
are given by the elementary symmetric functions
$e_j(x_1^r,\dots,x_n^r)$, $1 \le j \le n$ \cite{ATY1997}. Let $J'_+$
be the ideal of
$\mathcal{O}_X$ generated by these fundamental invariants and
$\Lambda' = \mathcal{O}_X/J'_+$ the quotient ring. It is also known
that the $G(r,n)$-module $\Lambda'$ is isomorphic to the group ring
$\mathbb{C}[G(r,n)]$, namely the left regular representation
\cite{ATY1997}. A
description of all irreducible components of $\Lambda'$ is known in
\cite{ATY1997}, in terms of what is called ``higher Specht
polynomials''. The irreducible representations of $G(r,n)$ are
parametrized by $r$-tuples of Young diagrams $(\lambda^1,\dots,\lambda^r)$
with $|\lambda^1|+\dots+|\lambda^r|=n$.

Now consider the restriction of the above action of $G(r,n)$ on
$\mathcal{O}_X$ to the subgroup $G(r,p,n)$. The fundamental invariants
are $e_j(x_1^r,\dots,x_n^r)$, $1 \le j \le n-1$, and
$e_n(x_1^d,\dots,x_n^d)$. Denote by $J_+$ the ideal generated by these
polynomials and let $\Lambda_0 = \mathcal{O}_X/J_+$. The representation
of $G(r,p,n)$ on $\Lambda_0$ is again isomorphic to the left regular
representation $\mathbb{C}[G(r,p,n)]$ \cite{MoritaYamada1998}.

\subsubsection{Higher Specht polynomials for $G(r,p,n)$}\label{subsubsec:shift}

Let $\mathcal{P}_{r,n}$ be the set of $r$-tuples of Young diagrams
$\lambda = (\lambda^1,\dots,\lambda^r)$ with
$|\lambda^1|+\dots+|\lambda^r|=n$. By filling each cell with a
positive integer so that every $j$ ($1 \le j \le n$) occurs once, we
obtain an $r$-tableau $T=(T^1,\dots,T^r)$ of shape $\lambda$. When the
number $k$ occurs in the component $T^i$, we write $k \in T^i$. The
set of $r$-tableaux of shape $\lambda$ is denoted $\mathrm{Tab}(\lambda)$.
An $r$-tableau $T=(T^1,\dots,T^r)$ is \emph{standard} if the numbers
are increasing along each column and each row of $T^\nu$
($1 \le \nu \le r$). The set of $r$-standard tableaux of shape
$\lambda$ is denoted $\mathrm{STab}(\lambda)$.

Let $S=(S^1,\dots,S^r) \in \mathrm{STab}(\lambda)$. We associate a
word $w(S)$ as follows: first we read each column of the component
$S^1$ from bottom to top, starting from the left; we continue this
procedure for $S^2$, and so on. For the word $w(S)$ we define the
index $i(w(S))$ inductively: the number $1$ has index $i(1)=0$; if the
number $k$ has index $i(k)=q$ and the number $k+1$ sits on the left
(resp.\ right) of $k$, then $k+1$ has index $q+1$ (resp.\ $q$).
Assigning the indices to the corresponding cells yields a shape
$\lambda$, each cell filled with a nonnegative integer, denoted
$i(S) = (i(S)^1,\dots,i(S)^r)$.

Let $T=(T^1,\dots,T^r)$ be an $r$-tableau of shape $\lambda$. For each
component $T^\nu$ ($1 \le \nu \le r$), the Young symmetrizer
$e_{T^\nu}$ of $T^\nu$ is defined by
\[
  e_{T^\nu} = \frac{1}{\alpha_{T^\nu}}
    \sum_{\sigma \in R(T^\nu)} \sum_{\tau \in C(T^\nu)} \mathrm{sgn}(\tau)\,\tau\sigma,
\]
where $\alpha_{T^\nu}$ is the product of the hook lengths of the shape
$\lambda^\nu$, and $R(T^\nu)$, $C(T^\nu)$ are the row- and
column-stabilizers of $T^\nu$, respectively.

We may regard a tableau $T$ on a Young diagram $\lambda$ as a map
$T : \{\text{cells of } \lambda\} \to \mathbb{Z}_{\ge 0}$ assigning to
a cell $\xi$ the value $T(\xi)$. For $S \in \mathrm{STab}(\lambda)$
and $T \in \mathrm{Tab}(\lambda)$, Ariki, Terasoma and Yamada
\cite{ATY1997} defined the higher Specht polynomial for $G(r,n)$ by
\[
  F_{S,T} = \prod_{\nu=1}^r \left( e_{T^\nu}\big(x^{r\, i(S)^\nu}_{T^\nu}\big)
    \prod_{k \in T^\nu} x_k^\nu \right),
  \qquad
  x^{r\, i(S)^\nu}_{T^\nu} = \prod_{\xi \in \lambda^\nu} x^{\,r\, i(S)^\nu(\xi)}_{T^\nu(\xi)}.
\]

Define the shift operator $\mathrm{Sh}$ on $\mathcal{P}_{r,n}$ (resp.\
on $\mathrm{Tab}(\lambda)$) by
\[
  \mathrm{Sh}(\lambda^1,\dots,\lambda^r) = (\lambda^r,\lambda^1,\dots,\lambda^{r-1})
  \qquad
  \big(\text{resp. } \mathrm{Sh}(T^1,\dots,T^r) = (T^r,T^1,\dots,T^{r-1})\big).
\]
For $\lambda,\mu \in \mathcal{P}_{r,n}$, define $\lambda \sim \mu$ if
$\mu = \mathrm{Sh}^{dj}(\lambda)$ for some $j=0,\dots,p-1$; this is an
equivalence relation on $\mathcal{P}_{r,n}$. Let
$b(\lambda) := |\{\mu \in \mathcal{P}_{r,n} \mid \lambda \sim \mu\}|$
and $e(\lambda) := p/b(\lambda)$. For $h=1,\dots,r$, let
\[
  \mathrm{STab}(\lambda)_h = \{ T=(T^1,\dots,T^r) \in \mathrm{STab}(\lambda) \mid
    1 \in T^\nu,\ 1 \le \nu \le h \}.
\]
For $\lambda \in \mathcal{P}_{r,n}$, fix $S \in \mathrm{STab}(\lambda)$
and $l=0,1,\dots,e(\lambda)-1$; for each $T \in \mathrm{STab}(\lambda)$
we define
\[
  F^l_{S,T} = \sum_{m=0}^{e(\lambda)-1} \xi^{lm\,d\,b(\lambda)}\,
    F_{S,\,\mathrm{Sh}^{mb(\lambda)}(T)}
  \;\in\; \Lambda_0.
\]
The family
\[
  \{ F^l_{S,T} \in \Lambda_0 \mid S \in \mathrm{STab}(\lambda)_d,\
     T \in \mathrm{STab}(\lambda)_{db(\lambda)},\
     l=0,1,\dots,e(\lambda)-1 \}
\]
is called the family of higher Specht polynomials for the complex
reflection group $G(r,p,n)$.

\begin{theorem}[{\cite{MoritaYamada1998}}]\label{thm:2.3}
Let $\lambda \in \mathcal{P}_{r,n}$, and for each
$S \in \mathrm{STab}(\lambda)_d$ and $0 \le l \le e(\lambda)-1$, put
\[
  V^l_S := \bigoplus_{T \in \mathrm{STab}(\lambda)_{db(\lambda)}} \mathbb{C} F^l_{S,T}
  \subset \Lambda_0.
\]
\begin{enumerate}[label=(\roman*)]
  \item The space $V^l_S(\lambda)$ affords an irreducible
    representation of the reflection group $G(r,p,n)$.
  \item We have the irreducible decomposition, as a representation of
    $G(r,p,n)$,
    \[
      \mathbb{C}[G(r,p,n)] = \bigoplus_{\lambda} \bigoplus_{S \in \mathrm{STab}(\lambda)_d}
        \bigoplus_{l=0}^{e(\lambda)-1} V^l_S(\lambda),
    \]
    where $\lambda$ runs over a system of representatives of
    $\mathcal{P}_{r,n}/\!\sim$.
  \item The higher Specht polynomials in
    \[
      \mathcal{F} = \{ F^l_{S,T} \mid S \in \mathrm{STab}(\lambda)_d,\
        T \in \mathrm{STab}(\lambda)_{db(\lambda)},\ \lambda \in \mathcal{P}_{r,n},\
        l=0,1,\dots,e(\lambda)-1 \}
    \]
    form a basis of the $\mathbb{C}[x_1,\dots,x_n]^{G(r,p,n)}$-module
    $\mathbb{C}[x_1,\dots,x_n]$.
\end{enumerate}
\end{theorem}


\section{Decomposition theorem for $G(r,p,n)$}\label{sec:general}

Let $W = G(r,p,n)$, $p \mid r$, $d := r/p$. Recall from
Subsection~\ref{subsec:specht} that $G(r,p,n)$ consists of
monomial matrices whose nonzero entries are $r$-th roots of unity
$\xi^{j}$, $0 \le j < r$, such that the $d$-th power of the product of
the nonzero entries equals $1$. Its reflections are of two kinds:

\begin{itemize}[leftmargin=1.8em]
 \item \emph{non-diagonal reflections}: for $1 \le i < j \le n$ and
 $0 \le k < r$, the map sending $x_i \mapsto \xi^{k} x_j$,
 $x_j \mapsto \xi^{-k} x_i$ and fixing the other coordinates;
 \item \emph{diagonal reflections} (present only when $d > 1$, i.e.\
 $p < r$): for $1 \le i \le n$ and $1 \le j < d$, the map
 $x_i \mapsto \xi^{jp} x_i$, fixing the other coordinates.
\end{itemize}

Recall $\Delta = \prod_{s \in S} \alpha_s(x)$ (Subsection~\ref{subsec:notation}).
Grouping this product by pair $\{i,j\}$ and by coordinate $i$: for
fixed $i<j$, the $r$ non-diagonal reflections ($k=0,\dots,r-1$)
contribute the linear forms $x_i - \xi^k x_j$, whose product is the
classical factorization
$\prod_{k=0}^{r-1}(x_i - \xi^k x_j) = x_i^r - x_j^r$; and for fixed
$i$, the $d-1$ diagonal reflections ($j=1,\dots,d-1$) each contribute
a factor proportional to $x_i$, together giving $x_i^{\,d-1}$.
Consequently the polynomial discriminant of $G(r,p,n)$ is
\[
 \Delta = \Delta_{G(r,p,n)}(x)
 \;=\; c_0 \prod_{i=1}^n x_i^{\,d-1} \prod_{1 \le i < j \le n} (x_i^r - x_j^r)
\]
for some nonzero constant $c_0$; the second product disintegrates to
$1$ when $d=1$ (i.e.\ $p=r$), recovering the case $W(D_n)$ of
Section~\ref{sec:Dn}.

We take as fundamental invariants
\[
 y_j = \sum_{i=1}^n x_i^{\,jr}, \quad 1 \le j \le n-1,
 \qquad
 y_n = (x_1 \cdots x_n)^d,
\]
and set $\mathcal{O}_Y = \mathbb{C}[x_1,\dots,x_n]^W = \mathbb{C}[y_1,\dots,y_n]$
(a polynomial ring on these $n$ algebraically independent invariants,
by the Chevalley--Shephard--Todd theorem \cite{Chevalley1955}),
\[
 \widetilde{\mathcal{O}}_X := \mathbb{C}[x_1,\dots,x_n,\Delta^{-1}], \quad
 \widetilde{\mathcal{O}}_Y := \mathbb{C}[y_1,\dots,y_n,\Delta^{-2}], \quad
 \widetilde{\mathcal{D}}_Y := \mathbb{C}\Big\langle y_1,\dots,y_n,
 \frac{\partial}{\partial y_1},\dots,\frac{\partial}{\partial y_n},
 \Delta^{-2}\Big\rangle .
\]

\subsection{The generalized Jacobian lemma}\label{subsec:jacobian}

\begin{lemma}\label{lem:jacobian-general}
Let $A$ be the transpose of the Jacobian matrix
$\big(\partial y_i / \partial x_j\big)_{1 \le i,j \le n}$. Then
\[
 \det A \;=\; (-1)^{n-1}\, d\, r^{\,n-1} (n-1)! \;
 \prod_{i=1}^n x_i^{\,d-1} \prod_{1 \le i < j \le n} (x_j^r - x_i^r).
\]
In particular $\det A = c \, \Delta_{G(r,p,n)}(x)$ for some nonzero
constant $c$, and $\widetilde{\mathcal{O}}_X$ is a $\widetilde{\mathcal{D}}_Y$-module.
\end{lemma}

\begin{proof}
For $1 \le j \le n-1$,
$\partial y_j / \partial x_i = j r \, x_i^{\,jr-1}$.
For $j=n$,
\[
 \frac{\partial y_n}{\partial x_i}
 = d\, x_i^{\,d-1} \prod_{k \neq i} x_k^{\,d}.
\]
Since $d \le r$, we have $jr - 1 \ge d-1$ for every $j \ge 1$, so the
$i$-th column of the Jacobian matrix has common factor $x_i^{\,d-1}$.
Writing $P := \prod_{k=1}^n x_k^{\,d}$, so that
$\prod_{k \neq i} x_k^{\,d} = P\, x_i^{-d}$, and setting
$u_i := x_i^{r}$, the $i$-th column becomes
\[
 x_i^{\,d-1} x_i^{-d}
 \big(r\, u_i,\; 2r\, u_i^2,\; \dots,\; (n-1) r\, u_i^{\,n-1},\; d P\big)^{T}
 = x_i^{-1} \big(r\, u_i,\; \dots,\; (n-1) r\, u_i^{\,n-1},\; d P\big)^{T}.
\]
Factor $x_i^{-1}$ out of every column, then $jr$ out of row $j$
($1 \le j \le n-1$) and $dP$ out of the last row:
\[
 \det A
 = \Big(\prod_{i=1}^n x_i^{-1}\Big) \cdot r^{\,n-1}(n-1)! \cdot dP
 \cdot \det N,
\]
where $N$ has rows $(u_1^{\,j}, \dots, u_n^{\,j})$, $1 \le j \le n-1$,
followed by the constant row $(1,\dots,1)$. Moving the last row to the
top costs $n-1$ transpositions, so by the classical Vandermonde
determinant,
\[
 \det N = (-1)^{n-1} \det \big(u_i^{\,j-1}\big)_{1 \le i,j \le n}
 = (-1)^{n-1} \prod_{i<j} (u_j - u_i).
\]
Since $\big(\prod_i x_i^{-1}\big) \cdot P = \prod_i x_i^{\,d-1}$ and
$u_j - u_i = x_j^r - x_i^r$, the formula follows. That
$\widetilde{\mathcal{O}}_X$ is a $\widetilde{\mathcal{D}}_Y$-module
now follows directly: since $\det A$ -- a nonzero scalar multiple of
$\Delta$ -- is itself nonzero on
$\widetilde{\mathcal{O}}_X$, the matrix $A$ is invertible there, and
$\partial/\partial y_1, \dots, \partial/\partial y_n$ are
recovered as $\mathbb{C}[x_1,\dots,x_n,\Delta^{-1}]$-linear
combinations of $\partial/\partial x_1, \dots, \partial/\partial x_n$.
\end{proof}

\begin{remark}
For $W(D_n) = G(2,2,n)$ one has $d=1$, and the formula reduces to
$\det A = (-1)^{n-1} 2^{n-1}(n-1)! \prod_{i<j}(x_j^2 - x_i^2)$; see
Section~\ref{sec:Dn}.
\end{remark}

\subsection{Simple components and their multiplicities}\label{subsec:simple-components}

The remainder of the argument is purely representation-theoretic,
and does not depend on the specific reflection group beyond
Lemma~\ref{lem:jacobian-general}.

\begin{proposition}\label{prop:embedding-DY}
There exists an injective map
$\mathbb{C}[W] \hookrightarrow \operatorname{Hom}_{\widetilde{\mathcal{D}}_Y}
(\widetilde{\mathcal{O}}_X, \widetilde{\mathcal{O}}_X)$.
\end{proposition}

The proof has two independent parts: injectivity of the natural map
$\mathbb{C}[W] \to \operatorname{End}_{\mathbb{C}}(\widetilde{\mathcal{O}}_X)$
(forgetting the $\widetilde{\mathcal{D}}_Y$-linearity condition), and
the fact that this map lands inside the centralizer of
$\widetilde{\mathcal{D}}_Y$, i.e.\ that every element of
$\mathbb{C}[W]$ commutes with $\widetilde{\mathcal{D}}_Y$.

\begin{lemma}[Injectivity]\label{lem:CW-injective}
The natural map $\mathbb{C}[W] \to \operatorname{End}_{\mathbb{C}}(\widetilde{\mathcal{O}}_X)$
is injective.
\end{lemma}

\begin{proof}
Recall from Subsection~\ref{subsec:specht} that
$\Lambda_0 = \mathcal{O}_X/J_+$ is isomorphic, as a
$\mathbb{C}[W]$-module, to the left regular representation
$\mathbb{C}[W]$. The regular representation is a \emph{faithful}
$\mathbb{C}[W]$-module: if $x \in \mathbb{C}[W]$ acts as zero on all of
$\mathbb{C}[W]$, then in particular $x = x \cdot 1 = 0$. Hence the
induced action of $\mathbb{C}[W]$ on the quotient module $\Lambda_0$
is already faithful. Since $\Lambda_0$ is a quotient of
$\mathcal{O}_X$ as a $\mathbb{C}[W]$-module, any
$x \in \mathbb{C}[W]$ annihilating all of $\mathcal{O}_X$ would also
annihilate $\Lambda_0$, forcing $x = 0$ by the previous sentence.
Thus $\mathbb{C}[W] \to \operatorname{End}_{\mathbb{C}}(\mathcal{O}_X)$
is injective, and localizing at $\Delta$ (which cannot introduce new
kernel elements, as $\widetilde{\mathcal{O}}_X \supset \mathcal{O}_X$)
gives injectivity into
$\operatorname{End}_{\mathbb{C}}(\widetilde{\mathcal{O}}_X)$.
\end{proof}

It remains to show that every element of $\mathbb{C}[W]$ commutes
with $\widetilde{\mathcal{D}}_Y$; combined with
Lemma~\ref{lem:CW-injective}, this proves
Proposition~\ref{prop:embedding-DY}. We give two independent proofs of
this commutation fact.

\begin{proof}[Proof of commutation with $\widetilde{\mathcal{D}}_Y$]
It suffices to show every element of $\mathbb{C}[W]$ commutes with
$y_1,\dots,y_n$ and with $\partial/\partial y_1, \dots, \partial/\partial
y_n$. The first is immediate since the $y_j$ are $W$-invariant. For
the second, let $K = \mathbb{C}(y_1,\dots,y_n)$ and
$L = \mathbb{C}(x_1,\dots,x_n)$. Then $K = L^{W}$ and $L/K$ is a
Galois extension with group $W$ -- this holds for any finite group
acting faithfully by field automorphisms on $L$, by Artin's theorem
\cite{Artin1942}, \cite[Chap.\ VI]{Lang2002},
regardless of the specific reflection group. Any derivation $D$ on $K$
extends uniquely to a derivation on $L$, and
$\sigma^{-1} D \sigma = D$ for every $\sigma \in W$, so $\sigma$
commutes with $D$.
\end{proof}

\begin{proof}[Alternative proof of commutation with $\widetilde{\mathcal{D}}_Y$]
We give a second, purely computational argument, which avoids
Artin's theorem and instead uses only the chain rule together with
the matrix $A$ of Lemma~\ref{lem:jacobian-general}.

For $w \in W$, let $M(w) \in GL_n(\mathbb{C})$ be the matrix of $w$ in
the coordinates $x_1,\dots,x_n$, so that $(wx)_k = \sum_i M(w)_{ki}x_i$,
and $W$ acts on functions by $(w\cdot f)(x) = f(w^{-1}x)$, as in
Subsection~\ref{subsec:specht}. Write $J(x) = \big(\partial y_i/\partial x_j\big)_{i,j}$,
so that $A = J^T$.

\emph{Step 1: an identity for $A$.} Since $y_j(wx) = y_j(x)$ for every
$x$, differentiating in $x_i$ via the chain rule gives
\[
  \sum_k \frac{\partial y_j}{\partial x_k}(wx)\, M(w)_{ki}
  = \frac{\partial y_j}{\partial x_i}(x),
\]
i.e.\ $J(wx)\,M(w) = J(x)$, equivalently
\[
  A(wx) = \big(M(w)^T\big)^{-1} A(x),
  \qquad\text{so}\qquad
  A(wx)^{-1} = A(x)^{-1} M(w)^T. \tag{$\star$}
\]

\emph{Step 2: commutation with $y_j$.} Immediate, since the $y_j$ are
$W$-invariant: $w\cdot(y_j f) = y_j\,(w\cdot f)$ for every $f$.

\emph{Step 3: commutation with $\partial/\partial y_j$.} Recall
$\partial/\partial y_j = \sum_i A^{-1}(x)_{ji}\,\partial/\partial x_i$.
For $f \in \widetilde{\mathcal{O}}_X$ and $x$ fixed,
\[
  \frac{\partial}{\partial y_j}(w\cdot f)(x)
  = \sum_i A^{-1}(x)_{ji}\,\frac{\partial}{\partial x_i}\big[f(w^{-1}x)\big]
  = \sum_{i,k} A^{-1}(x)_{ji}\, M(w^{-1})_{ki}\,
    \frac{\partial f}{\partial x_k}(w^{-1}x)
  = \sum_k \big[A^{-1}(x) M(w^{-1})^T\big]_{jk}\,
    \frac{\partial f}{\partial x_k}(w^{-1}x).
\]
Set $y := w^{-1}x$. Applying $(\star)$ at $y$ gives
$A(x)^{-1} = A(wy)^{-1} = A(y)^{-1}M(w)^T$, hence
\[
  A^{-1}(x)\,M(w^{-1})^T
  = A(y)^{-1}M(w)^T M(w^{-1})^T
  = A(y)^{-1}\big(M(w^{-1})M(w)\big)^T
  = A(y)^{-1} = A^{-1}(w^{-1}x).
\]
Therefore
\[
  \frac{\partial}{\partial y_j}(w\cdot f)(x)
  = \sum_k A^{-1}(w^{-1}x)_{jk}\, \frac{\partial f}{\partial x_k}(w^{-1}x)
  = \left(\frac{\partial f}{\partial y_j}\right)(w^{-1}x)
  = w\cdot\!\left(\frac{\partial f}{\partial y_j}\right)(x),
\]
so $w$ commutes with $\partial/\partial y_j$ for every $j$, and hence
with all of $\widetilde{\mathcal{D}}_Y$.
\end{proof}

\begin{remark}
This second proof has the advantage of relying on no external
citation -- it reuses directly the matrix $A$ already introduced in
Lemma~\ref{lem:jacobian-general} -- at the cost of a longer,
index-chasing computation, compared to the conceptual brevity of the
Galois-theoretic argument above.
\end{remark}

\begin{corollary}[The group algebra as the full centralizer]\label{cor:iso-CW}
\[
  \mathbb{C}[W] \;\cong\; \operatorname{Hom}_{\widetilde{\mathcal{D}}_Y}
  \big(\widetilde{\mathcal{O}}_X, \widetilde{\mathcal{O}}_X\big).
\]
\end{corollary}

\begin{proof}
By Proposition~\ref{prop:embedding-DY}, the natural map is injective;
we prove surjectivity. Write $L = \operatorname{Frac}(\widetilde{\mathcal O}_X)$
and $K = \operatorname{Frac}(\widetilde{\mathcal O}_Y)$.

\smallskip\noindent\emph{Step 1: reduction to $L$.} Let
$\phi \in \operatorname{Hom}_{\widetilde{\mathcal D}_Y}(\widetilde{\mathcal O}_X,\widetilde{\mathcal O}_X)$.
Since $\phi$ commutes with multiplication by $y_1,\dots,y_n$, it is
$\widetilde{\mathcal O}_Y$-linear, and therefore extends uniquely to a
$K$-linear endomorphism of $L$ -- still denoted $\phi$ -- via
$\phi(f/g) := \phi(f)/g$ for $f \in \widetilde{\mathcal O}_X$,
$g \in \widetilde{\mathcal O}_Y \setminus \{0\}$. By the quotient rule,
$\phi$ continues to commute with the extended derivations
$\widetilde D_1,\dots,\widetilde D_n$ of $L$
(Lemma~\ref{lem:der-span}).\footnote{Lemma~\ref{lem:der-span} and
Lemma~\ref{lem:rigidity}, invoked later in this proof, are established
independently below and do not rely on the present corollary, so no
circularity arises from citing them here.}

\smallskip\noindent\emph{Step 2: $\operatorname{End}_K(L)$ as an
$L[W]$-module.} By Dedekind's lemma on the linear independence of
characters \cite[Chap.\ VI]{Lang2002}, the elements of $W$, viewed as
$K$-linear endomorphisms of $L$, are linearly independent over $L$
(where $L$ acts by left multiplication,
$a \cdot \Phi := (x \mapsto a\,\Phi(x))$ for $a \in L$). Since
\[
  \dim_K \operatorname{End}_K(L) = [L:K]^2 = |W|^2 = |W| \cdot \dim_K L,
\]
this linear independence forces, by a dimension count,
\[
  \operatorname{End}_K(L) = \bigoplus_{w \in W} L \cdot w,
\]
i.e.\ every $\Phi \in \operatorname{End}_K(L)$ is \emph{uniquely} of
the form $\Phi = \sum_{w \in W} a_w\, w$ for functions $a_w \in L$.

\smallskip\noindent\emph{Step 3: flatness forces constant
coefficients.} Write $\phi = \sum_w a_w\, w$ as in Step 2, and fix
$j \in \{1,\dots,n\}$. Since $\phi$ commutes with $\widetilde D_j$,
and each $w \in W$ already commutes with $\widetilde D_j$
(Proposition~\ref{prop:embedding-DY}), the Leibniz rule gives, for
every $x \in L$,
\[
  0 = \widetilde D_j(\phi(x)) - \phi(\widetilde D_j x)
    = \sum_{w \in W} \widetilde D_j(a_w)\, w(x).
\]
As this holds for all $x \in L$, Dedekind's lemma forces
$\widetilde D_j(a_w) = 0$ for every $j$ and every $w \in W$. By
Lemma~\ref{lem:der-span}, the derivations
$\widetilde D_1,\dots,\widetilde D_n$ span $\operatorname{Der}_{\mathbb{C}}(L)$
over $L$, so each $a_w$ is annihilated by \emph{every}
$\mathbb{C}$-derivation of $L$; as in the proof of
Lemma~\ref{lem:rigidity}, this forces $a_w \in \mathbb{C}$, since an
element of $L$ algebraic over $\mathbb{C}$ must already lie in
$\mathbb{C}$, $\mathbb{C}$ being algebraically closed.

Hence $\phi = \sum_{w \in W} a_w\, w$ with all $a_w \in \mathbb{C}$,
exhibiting $\phi$ as the restriction to $\widetilde{\mathcal O}_X$ of
an element of $\mathbb{C}[W]$. This proves surjectivity, and the
corollary follows.
\end{proof}

By Maschke's theorem \cite[Chap.\ XVIII]{Lang2002}, $\mathbb{C}[W]$
is semisimple, and by Theorem~\ref{thm:2.3},
\[
 \mathbb{C}[W] = \bigoplus_{\lambda \in \mathcal{P}_{r,n}/\sim}
 \bigoplus_{l=0}^{e(\lambda)-1} R^l_\lambda, \qquad
 R^l_\lambda = \bigoplus_{S \in \mathrm{STab}(\lambda)_d} V^l_S(\lambda),
\]
with primitive central idempotents $r^l_\lambda$, $\lambda \in
\mathcal{P}_{r,n}/\sim$, $l = 0,\dots,e(\lambda)-1$.

\begin{example}[Block decomposition of $\mathbb{C}[G(2,1,2)$]\label{ex:CW-decomp}
Let $W = G(2,1,2)$ (the dihedral group of order $8$; $r=2$, $p=1$,
$d=2$). Since $p=1$, $e(\lambda)=1$ for every
$\lambda \in \mathcal{P}_{2,2}$ (no further splitting by $l$), and
$\mathcal{P}_{2,2}$ has five elements: four singleton shapes (each
giving $f_\lambda=1$) and one ``pair'' shape $\lambda_0$ (giving
$f_{\lambda_0}=2$) -- these are worked out in detail, together with
explicit generating polynomials, in Subsection~\ref{subsec:Gr1n}. The
block decomposition displayed above reads concretely
\[
  \mathbb{C}[W] = \underbrace{\mathbb{C}}_{\mathrm{triv}} \oplus
    \underbrace{\mathbb{C}}_{\mathrm{sign}} \oplus
    \underbrace{\mathbb{C}}_{\chi^0} \oplus
    \underbrace{\mathbb{C}}_{\chi^1} \oplus
    \underbrace{\mathrm{Mat}_2(\mathbb{C})}_{R_{\lambda_0}},
\]
of total dimension $1+1+1+1+4 = 8 = |W|$: the four singleton shapes
give four $1\times 1$ blocks $R_\lambda \cong \mathbb{C}$ (the four
linear characters of $W$), while $\lambda_0$ gives the single
$2\times 2$ block $R_{\lambda_0} \cong \operatorname{Mat}_2(\mathbb{C})$
corresponding to the unique $2$-dimensional irreducible representation
of $W$. The primitive central idempotent of each $1\times 1$ block is
simply the character-averaging projector
$r_\lambda = \frac{1}{8}\sum_{g \in W} \overline{\chi_\lambda(g)}\, g$
(the characters here happen to be real-valued, so the bar is
inessential, but we keep it for consistency with the general formula
below); the block
$R_{\lambda_0}$ splits further into two primitive (non-central)
idempotents $e^{\lambda_0}_1, e^{\lambda_0}_2$, computed explicitly in
Subsection~\ref{subsec:Gr1n-idem}.
\end{example}

\begin{lemma}[Double centralizer lemma]\label{lem:double-centralizer}
Let $R$ be a $\mathbb{C}$-algebra and $A$ a finite-dimensional
semisimple $\mathbb{C}$-algebra ($\mathbb{C}=\overline{\mathbb{C}}$)
acting on a left $R$-module $M$ by $R$-endomorphisms, such that the
induced map $A \to \operatorname{End}_R(M)$ is an \emph{isomorphism}.
Write $A = \bigoplus_\lambda \operatorname{Mat}_{f_\lambda}(\mathbb{C})$
(Wedderburn's structure theorem, \cite[Chap.\ II, \S4]{Lang2002}),
with primitive idempotents
$e_1^\lambda,\dots,e_{f_\lambda}^\lambda$ in each block. Then:
\begin{enumerate}[label=(\alph*)]
  \item each $e_i^\lambda M$ is a simple $R$-module;
  \item $e_i^\lambda M \cong e_j^\lambda M$ as $R$-modules for all
    $i,j$ in the same block $\lambda$;
  \item $M = \bigoplus_{\lambda,i} e_i^\lambda M$, and
    $N_\lambda := e_1^\lambda M$ are pairwise non-isomorphic for
    $\lambda \ne \mu$.
\end{enumerate}
\end{lemma}

\begin{proof}
Since $A = \operatorname{End}_R(M)$ exactly, for any idempotents
$e_i^\lambda, e_j^\mu$,
\[
  \operatorname{Hom}_R(e_i^\lambda M,\, e_j^\mu M) \;\cong\; e_i^\lambda A e_j^\mu,
\]
via $\varphi \mapsto e_i^\lambda \varphi(-) e_j^\mu$, which is
bijective precisely because $A$ is \emph{all} of
$\operatorname{End}_R(M)$. If $\lambda \ne \mu$, then
$e_i^\lambda A e_j^\mu = 0$ (orthogonal blocks), giving the
non-isomorphism in (c). If $\lambda = \mu$, then
$e_i^\lambda A e_j^\lambda \cong \mathbb{C}$ (the $(i,j)$ entry of a
matrix algebra); in particular
$\operatorname{End}_R(e_i^\lambda M) \cong \mathbb{C}$, so by Schur's
lemma \cite[\S 1.4]{FultonHarris1991} $e_i^\lambda M$ is simple,
giving (a); and
$\operatorname{Hom}_R(e_i^\lambda M, e_j^\lambda M) \cong \mathbb{C} \ne 0$
gives the isomorphism in (b), again by Schur's lemma.
\end{proof}

\paragraph{Idempotents of $\mathbb{C}[G(r,p,n)]$.}
Before stating Theorem~\ref{thm:main-general}, we make explicit the
central and primitive idempotents to which it applies, in terms of
the character theory of $W = G(r,p,n)$ furnished by
Theorem~\ref{thm:2.3}.

For $\lambda \in \mathcal{P}_{r,n}/\sim$ and $0 \le l < e(\lambda)$,
let $\chi^l_\lambda$ denote the character of the irreducible
representation $V^l_S(\lambda)$ -- independent of the choice of
$S \in \mathrm{STab}(\lambda)_d$, since all such $V^l_S(\lambda)$ are
isomorphic (Theorem~\ref{thm:2.3}(i)) -- of dimension
$f^l_\lambda := \dim_{\mathbb{C}} V^l_S(\lambda)$. By the standard
orthogonality relations for characters of a finite group
\cite[\S 2.3]{FultonHarris1991}, the corresponding primitive central
idempotent of $\mathbb{C}[W]$ is
\[
  r^l_\lambda = \frac{f^l_\lambda}{|W|} \sum_{g \in W} \overline{\chi^l_\lambda(g)}\; g,
\]
satisfying $\sum_{\lambda,l} r^l_\lambda = 1$ and
$r^l_\lambda r^j_\mu = 0$ for $(\lambda,l) \ne (\mu,j)$; this is
exactly the identity element of the block $R^l_\lambda$ displayed
above.

Within each block, fix a matrix realization
$\rho^l_\lambda : W \to \mathrm{GL}_{f^l_\lambda}(\mathbb{C})$ of
$V^l_S(\lambda)$ for one choice of $S$. The matrix-coefficient
formula \cite[\S 2.4]{FultonHarris1991} gives, for
$1 \le i,j \le f^l_\lambda$,
\[
  e^l_{\lambda,ij} = \frac{f^l_\lambda}{|W|} \sum_{g \in W}
    \overline{\rho^l_\lambda(g)_{ji}}\; g \;\in\; R^l_\lambda \subset \mathbb{C}[W],
\]
satisfying the matrix-unit relations
$e^l_{\lambda,ij}\, e^l_{\lambda,kl'} = \delta_{jk}\, e^l_{\lambda,il'}$
and $\sum_i e^l_{\lambda,ii} = r^l_\lambda$. In particular, each
diagonal element $e^l_{\lambda,ii}$ ($1 \le i \le f^l_\lambda$) is a
\emph{primitive idempotent} of $\mathbb{C}[W]$; relabelling
$e^{\lambda,l}_i := e^l_{\lambda,ii}$, these are exactly the
idempotents $e_1^\lambda,\dots,e_{f_\lambda}^\lambda$ of
Lemma~\ref{lem:double-centralizer} for the block $(\lambda,l)$. Every
primitive idempotent of $\mathbb{C}[W]$ arises this way, for a unique
pair $(\lambda,l)$ and some $i$ (up to the residual freedom of
choosing the basis diagonalizing $\rho^l_\lambda$). It is this
explicit family $\{e^{\lambda,l}_i\}$ to which
Theorem~\ref{thm:main-general} below applies.

\begin{theorem}\label{thm:main-general}
Fix $\lambda \in \mathcal{P}_{r,n}/\sim$ and $0 \le l < e(\lambda)$,
and enumerate $\mathrm{STab}(\lambda)_d = \{S_1,\dots,S_{f^l_\lambda}\}$.
For every $i \in \{1,\dots,f^l_\lambda\}$, writing $e^{\lambda,l}_i$
for the corresponding primitive idempotent of $\mathbb{C}[W]$
(as described above):
\begin{enumerate}[label=(\roman*)]
 \item $e^{\lambda,l}_i\,\widetilde{\mathcal{O}}_X$ is a nontrivial
 $\widetilde{\mathcal{D}}_Y$-submodule of $\widetilde{\mathcal{O}}_X$;
 \item the $\widetilde{\mathcal{D}}_Y$-module
 $e^{\lambda,l}_i\,\widetilde{\mathcal{O}}_X$ is simple;
 \item there exists a higher Specht polynomial $F^l_{S_i,T}$, for some
 $T \in \mathrm{STab}(\lambda)_{d\,b(\lambda)}$, such that
 $e^{\lambda,l}_i\,\widetilde{\mathcal{O}}_X = \widetilde{\mathcal{D}}_Y \, F^l_{S_i,T}$.
\end{enumerate}
\end{theorem}

\begin{proof}
We apply Lemma~\ref{lem:double-centralizer} with $R = \widetilde{\mathcal D}_Y$,
$M = \widetilde{\mathcal O}_X$, $A = \mathbb{C}[W]$. The hypotheses hold
precisely by what we have already established: $A$ is semisimple
(Maschke), and $A \to \operatorname{End}_R(M)$ is an isomorphism by
Corollary~\ref{cor:iso-CW}. Applying the lemma to the block
$(\lambda,l)$ and idempotent $e^{\lambda,l}_i$: part~(c) of the lemma
exhibits $e^{\lambda,l}_i\,\widetilde{\mathcal{O}}_X$ as a genuine
direct-summand submodule of $\widetilde{\mathcal{O}}_X$, while
part~(a) shows it is simple -- and simple modules are by definition
nonzero, so part~(a) alone already gives both nontriviality and
simplicity. Together, (a) and (c) give parts~(i) and~(ii) of the
theorem, with no case-by-case verification required.

For~(iii): this tableau-indexing of the primitive idempotents is not
an extra choice grafted onto Lemma~\ref{lem:double-centralizer} -- it
is exactly Theorem~\ref{thm:2.3}(ii)'s decomposition
$\mathbb{C}[W] = \bigoplus_{\lambda,l} \bigoplus_{S \in \mathrm{STab}(\lambda)_d} V^l_S(\lambda)$,
whose summands are already indexed by tableaux $S$; enumerating
$\mathrm{STab}(\lambda)_d = \{S_1,\dots,S_{f^l_\lambda}\}$ and matching
each $V^l_{S_i}(\lambda)$ with the corresponding minimal left ideal
$\mathbb{C}[W]e^{\lambda,l}_i$ recovers exactly the matrix units
$e^l_{\lambda,ii}$ constructed above, i.e.\ $e^{\lambda,l}_i$
corresponds precisely to the tableau $S_i$. By
Theorem~\ref{thm:2.3}(ii)--(iii), the associated block
$V^l_{S_i}(\lambda) = \mathbb{C}[W]F^l_{S_i,T}$ is cyclic, generated by
a higher Specht polynomial $F^l_{S_i,T}$. Hence
$F^l_{S_i,T} \in \mathbb{C}[W]F^l_{S_i,T} \subset e^{\lambda,l}_i\,\widetilde{\mathcal O}_X$
is a nonzero element of $e^{\lambda,l}_i\,\widetilde{\mathcal O}_X$; by
the simplicity already established in~(ii), any nonzero element
generates the whole module, so
$e^{\lambda,l}_i\,\widetilde{\mathcal O}_X = \widetilde{\mathcal D}_Y\, F^l_{S_i,T}$.
\end{proof}

\begin{corollary}\label{cor:iso-simple}
With the above notation, for fixed $\lambda \in \mathcal{P}_{r,n}/\sim$
and $0 \le l < e(\lambda)$,
\[
  e^{\lambda,l}_i\,\widetilde{\mathcal{O}}_X \;\cong_{\widetilde{\mathcal{D}}_Y}\;
  e^{\lambda,l}_j\,\widetilde{\mathcal{O}}_X
  \qquad \text{for all } 1 \le i,j \le f^l_\lambda.
\]
\end{corollary}

\begin{proposition}\label{prop:decomposition-general}
With the enumeration $\mathrm{STab}(\lambda)_d = \{S_1,\dots,S_{f^l_\lambda}\}$
of Theorem~\ref{thm:main-general}, write
$F^l_{\lambda,i} := F^l_{S_i,T}$ for the generator of
$e^{\lambda,l}_i\,\widetilde{\mathcal O}_X$ given by
Theorem~\ref{thm:main-general}(iii). Then:
\begin{enumerate}[label=(\roman*)]
\item
\[
 \widetilde{\mathcal{O}}_X
 = \bigoplus_{\lambda \in \mathcal{P}_{r,n}/\sim} \bigoplus_{l=0}^{e(\lambda)-1}
 \bigoplus_{i=1}^{f^l_\lambda} \widetilde{\mathcal{D}}_Y F^l_{\lambda,i} ;
\]
\item fixing, for each $(\lambda,l)$, the representative $i=1$,
\[
 \widetilde{\mathcal{O}}_X
 = \bigoplus_{\lambda \in \mathcal{P}_{r,n}/\sim} \bigoplus_{l=0}^{e(\lambda)-1}
 f^l_\lambda \, \widetilde{\mathcal{D}}_Y F^l_{\lambda,1}.
\]
\end{enumerate}
\end{proposition}

\begin{proof}
Same argument as in the classical case, using
Theorem~\ref{thm:main-general} and Corollary~\ref{cor:iso-simple} in
place of an ad hoc argument for a specific reflection group: part (i)
is $\widetilde{\mathcal O}_X = \bigoplus_{\lambda,l,i} e^{\lambda,l}_i \widetilde{\mathcal O}_X$
(Lemma~\ref{lem:double-centralizer}(c)) rewritten via
Theorem~\ref{thm:main-general}(iii); part (ii) groups the
$f^l_\lambda$ isomorphic summands $e^{\lambda,l}_i\widetilde{\mathcal O}_X$
($i=1,\dots,f^l_\lambda$) of a fixed block, using
Corollary~\ref{cor:iso-simple}, into $f^l_\lambda$ copies of the single
representative $\widetilde{\mathcal D}_Y F^l_{\lambda,1}$.
\end{proof}

\begin{example}[Proposition~\ref{prop:decomposition-general} for $B_2 = G(2,1,2)$]\label{ex:prop29}
We illustrate both parts of the proposition on the group
$W = G(2,1,2)$ (so $r=2$, $p=1$, $d=2$, $|W|=8$), using the explicit
higher Specht polynomials of Subsection~\ref{subsec:Gr1n} (case $r=2$).
Here $\mathcal{P}_{2,2}$ has $5$ elements, with $e(\lambda)=1$
throughout since $p=1$ (so we omit the superscript $l=0$ below): two
singleton shapes $(2)$ (components $\nu=0,1$), two singleton shapes
$(1,1)$ (components $\nu=0,1$), and one ``pair'' shape $\lambda_0$
(single boxes in components $0$ and $1$), the last with
$f_{\lambda_0} = 2$ and $\mathrm{STab}(\lambda_0)_d = \{S_1,S_2\}$
(writing $S_1=T_a$, $S_2=T_b$ in the notation of
Subsection~\ref{subsec:Gr1n}). The five standard tableaux are, in
order:
\[
  \nu{=}0:\begin{array}{|c|c|}\hline 1&2\\\hline\end{array}\ ,\quad
  \nu{=}1:\begin{array}{|c|c|}\hline 1&2\\\hline\end{array}\ ,\quad
  \nu{=}0:\begin{array}{|c|}\hline 1\\\hline 2\\\hline\end{array}\ ,\quad
  \nu{=}1:\begin{array}{|c|}\hline 1\\\hline 2\\\hline\end{array}\ ,
\]
\[
  T_a = \Big(0:\begin{array}{|c|}\hline 1\\\hline\end{array},\
    1:\begin{array}{|c|}\hline 2\\\hline\end{array}\Big),
  \qquad
  T_b = \Big(0:\begin{array}{|c|}\hline 2\\\hline\end{array},\
    1:\begin{array}{|c|}\hline 1\\\hline\end{array}\Big),
\]
generating $F_0,F_1,F'_0,F'_1$ and (together with the corresponding
$T\in\{T_a,T_b\}$) the two pairs of generators for $\lambda_0$,
respectively.

\smallskip\noindent\emph{Part (i).} Summing over all $\lambda$ and all
$i=1,\dots,f_\lambda$ gives \emph{six} summands (one each for the four
singleton shapes, two for $\lambda_0$, matching
$\sum_\lambda f_\lambda = 1+1+1+1+2=6$):
\[
  \widetilde{\mathcal O}_X =
  \widetilde{\mathcal D}_Y\cdot 1 \;\oplus\;
  \widetilde{\mathcal D}_Y\cdot x_1x_2 \;\oplus\;
  \widetilde{\mathcal D}_Y\cdot(x_1^2-x_2^2) \;\oplus\;
  \widetilde{\mathcal D}_Y\cdot x_1x_2(x_1^2-x_2^2) \;\oplus\;
  \widetilde{\mathcal D}_Y\cdot x_2 \;\oplus\;
  \widetilde{\mathcal D}_Y\cdot x_1^2x_2,
\]
where the last two summands are
$\widetilde{\mathcal D}_Y F_{\lambda_0,1}$ and
$\widetilde{\mathcal D}_Y F_{\lambda_0,2}$ (taking, for each, either
polynomial in the corresponding pair of Subsection~\ref{subsec:Gr1n}
as a generator -- by Theorem~\ref{thm:main-general}(ii) the module
$e_1^{\lambda_0}\widetilde{\mathcal O}_X$ is simple, so $x_1$ and $x_2$
generate the \emph{same} rank-$2$ submodule
$\widetilde{\mathcal D}_Y\cdot x_2 = \widetilde{\mathcal D}_Y\cdot x_1$).
As a check, the $\widetilde{\mathcal O}_Y$-ranks add up to
$1+1+1+1+2+2 = 8 = |W|$.

\smallskip\noindent\emph{Part (ii).} Fixing $i=1$ for $\lambda_0$ (so
$f_{\lambda_0}=\dim V(\lambda_0)=2$), the two isomorphic summands of
part (i) coming from $i=1$ and $i=2$
(Corollary~\ref{cor:iso-simple}: $e_1^{\lambda_0}\widetilde{\mathcal O}_X \cong
e_2^{\lambda_0}\widetilde{\mathcal O}_X$) collapse into a single
isomorphism type with multiplicity $2$:
\[
  \widetilde{\mathcal O}_X =
  \widetilde{\mathcal D}_Y\cdot 1 \;\oplus\;
  \widetilde{\mathcal D}_Y\cdot x_1x_2 \;\oplus\;
  \widetilde{\mathcal D}_Y\cdot(x_1^2-x_2^2) \;\oplus\;
  \widetilde{\mathcal D}_Y\cdot x_1x_2(x_1^2-x_2^2) \;\oplus\;
  \big(\widetilde{\mathcal D}_Y\cdot x_2\big)^{\oplus 2}.
\]
The rank count is unchanged: $1+1+1+1+2\cdot 2 = 8$. This is exactly
the grouping by central idempotent: the first four summands correspond
to the four $1$-dimensional characters of $B_2$ (a dihedral group of
order $8$), while the last, doubled summand corresponds to its unique
$2$-dimensional irreducible representation.
\end{example}

\begin{remark}
Taking $r=p=2$ recovers the decomposition theorem for $W(D_n)$,
worked out in detail in Section~\ref{sec:Dn} below. Taking $p=r$
(so $d=1$) gives the decomposition for the index-$r$ subgroup
$G(r,r,n)$; taking $p=1$ (so $d=r$) gives it for the full monomial
group $G(r,1,n)$. Both cases are treated explicitly in
Section~\ref{sec:other-cases}.
\end{remark}

\subsection{Categorical equivalence and new consequences}\label{subsec:galois-descent}

The abstract Galois-descent equivalence of categories used in this
subsection -- for \emph{any} finite group acting faithfully on a
field -- is due to
Nonkan\'e \cite[\S 2.4]{Nonkane2019}, where it is applied to $L/K$ for
a single symmetric group $W = \mathcal{S}_n$. It is recalled
verbatim -- explicitly citing \cite{Nonkane2019} for it -- and applied
to products of symmetric groups by
Nonkan\'e--Todjihounde \cite[Prop.\ 3.10]{NonkaneTodjihounde2023}, for
$L/K$ with $K = \operatorname{Frac}(\mathcal{O}_X)^{\mathcal{S}_{n_1}\times\cdots\times\mathcal{S}_{n_r}}$.
What is new here is twofold: the
application of this equivalence to $W = G(r,p,n)$ (Theorem~\ref{thm:descent}
below, stated in our notation but not reproved, since a proof is
already given in \cite{Nonkane2019}), and the two consequences of
Corollaries~\ref{cor:twisted-invariants} and~\ref{cor:characterization},
which do not appear in \cite{Nonkane2019,NonkaneTodjihounde2023}. We
also record, in (F1)--(F3) and Lemmas~\ref{lem:descent-iso}--\ref{lem:rigidity}
below, the elementary Galois-theoretic facts underlying the
equivalence; these are classical (Artin's theorem, the Normal Basis
Theorem, and Kähler differentials) rather than specific to
\cite{Nonkane2019,NonkaneTodjihounde2023}, and we state them here
because Corollaries~\ref{cor:twisted-invariants}
and~\ref{cor:characterization} -- our own new results -- use them
directly.

Let $K = \operatorname{Frac}(\mathcal{O}_Y)$ and
$L = \operatorname{Frac}(\mathcal{O}_X)$. By Proposition~\ref{prop:embedding-DY}
(either proof), $K = L^W$ and $L/K$ is Galois with group $W = G(r,p,n)$.

\begin{theorem}[Galois descent, {\cite[\S 2.4]{Nonkane2019}}]\label{thm:descent}
Let $\operatorname{Rep}(W)$ denote the category of finite-dimensional
$\mathbb{C}$-representations of $W$. The functor
\[
  \nabla : \operatorname{Rep}(W) \longrightarrow \widetilde{\mathcal{D}}_Y\text{-}\mathrm{Mod},
  \qquad V \longmapsto \big(\widetilde{\mathcal{O}}_X \otimes_{\mathbb{C}} V\big)^{W}
\]
(with $W$ acting diagonally) is fully faithful, hence defines an
equivalence of categories between $\operatorname{Rep}(W)$ and the full
subcategory of $\widetilde{\mathcal{D}}_Y$-modules that become trivial
after the base change
$\widetilde{\mathcal{O}}_X \otimes_{\widetilde{\mathcal{O}}_Y} (-)$ --
that is, isomorphic to $\widetilde{\mathcal{O}}_X \otimes_{\mathbb{C}} U$
for some finite-dimensional $\mathbb{C}$-vector space $U$, with $W$
acting only through its own action on the $\widetilde{\mathcal{O}}_X$
factor (as opposed to the diagonal action, which also permutes the
$U$-coordinates) -- the essential image of $\nabla$ by definition. A
quasi-inverse is
$\operatorname{Loc}(M) = \big(\widetilde{\mathcal{O}}_X \otimes_{\widetilde{\mathcal{O}}_Y} M\big)^{\widetilde{\mathcal{D}}_Y\text{-flat part}}$,
i.e.\ the $W$-module of flat sections of the base-changed connection.
This instance, for $W = G(r,p,n)$, is new; the underlying equivalence
itself, for a general finite group, is due to
\cite[\S 2.4]{Nonkane2019} (see also
\cite[Prop.\ 3.10]{NonkaneTodjihounde2023}, which recalls it for
products of symmetric groups), and we do not reprove it
here.
\end{theorem}

Write $L = \operatorname{Frac}(\widetilde{\mathcal{O}}_X)$,
$K = \operatorname{Frac}(\widetilde{\mathcal{O}}_Y)$ (the same fields
as above, since localizing at $\Delta$ does not change the fraction
field). The facts and lemmas below are recorded because
Corollaries~\ref{cor:twisted-invariants} and~\ref{cor:characterization}
use them directly; the reader willing to take
Theorem~\ref{thm:descent} as given may skip ahead to
Corollary~\ref{cor:twisted-invariants}. We
use three classical facts:

\begin{itemize}
  \item[(F1)] \emph{Extension of derivations.} Every derivation of $K$
    extends uniquely to a derivation of $L$ (proved in the first
    proof of Proposition~\ref{prop:embedding-DY} above), and this
    extension commutes with the $W$-action.
  \item[(F2)] \emph{Degree of the extension.} $[L:K] = |W|$ (Artin's
    theorem on fixed fields of finite group actions, see
    \cite{Artin1942} or e.g.\
    \cite[Chap.\ VI]{Lang2002}).
  \item[(F3)] \emph{Normal Basis Theorem} (see e.g.\ \cite[Chap.\ VI]{Lang2002}).
    $L \cong K[W]$ as $K[W]$-modules, where $W$ acts on $K[W]$ by left
    translation.
\end{itemize}

\begin{lemma}[Descent isomorphism]\label{lem:descent-iso}
For every finite-dimensional $\mathbb{C}$-vector space $V$ with a
$W$-action, the multiplication map
\[
  \mu_V : L \otimes_K \big(L \otimes_{\mathbb{C}} V\big)^{W} \longrightarrow L \otimes_{\mathbb{C}} V,
  \qquad a \otimes m \longmapsto am
\]
is an isomorphism of $L$-vector spaces.
\end{lemma}

\begin{proof}
\emph{Injectivity.} Suppose $\sum_{i=1}^r a_i \otimes m_i \mapsto 0$
with $m_1,\dots,m_r \in (L\otimes V)^W$ linearly independent over $K$
and $r$ minimal among nonzero relations; then all $a_i \ne 0$, and
after scaling we may assume $a_1 = 1$. For $\sigma \in W$, applying
$\sigma$ to $\sum a_i m_i = 0$ and using $\sigma(m_i) = m_i$ gives
$\sum \sigma(a_i) m_i = 0$; subtracting the original relation gives
$\sum_{i \ge 2} (\sigma(a_i) - a_i) m_i = 0$, a shorter relation
(the $i=1$ term vanishes since $\sigma(a_1)-a_1=1-1=0$), so by
minimality $\sigma(a_i) = a_i$ for all $i,\sigma$, i.e.\ all
$a_i \in L^W = K$. But then $\sum a_i m_i = 0$ is a $K$-linear
relation among the $K$-linearly independent $m_i$, forcing all
$a_i = 0$ -- contradicting $a_1 = 1 \ne 0$. Hence $\mu_V$ is
injective.

\emph{Dimension count.} By (F3), $L \cong K[W]$ as $K[W]$-modules
(where $W$ acts on $L$ via its Galois action, matching the left
regular representation on $K[W]$). Under this identification, the map
$g \otimes v \mapsto g \otimes g^{-1}v$ gives a $K[W]$-module
isomorphism $K[W] \otimes_{\mathbb{C}} V \xrightarrow{\sim} K[W] \otimes_{\mathbb{C}} V$
from the diagonal action to the action on the first factor only (a
direct check, valid for any group and any base ring): indeed, for
$h \in W$, the diagonal action sends
$g \otimes v \mapsto hg \otimes hv$, which under the map goes to
$hg \otimes (hg)^{-1}hv = hg \otimes g^{-1}v$, matching the action of
$h$ on $g \otimes g^{-1}v$ when $W$ acts only on the first factor.
Hence, as $K[W]$-modules,
$L \otimes_{\mathbb{C}} V \cong K[W]^{\oplus \dim_{\mathbb{C}} V}$ with
$W$ acting only on the $K[W]$ factors, so
\[
  \dim_K (L \otimes_{\mathbb{C}} V)^W
  = \dim_{\mathbb{C}} V \cdot \dim_K \big(K[W]\big)^W
  = \dim_{\mathbb{C}} V,
\]
since $(K[W])^W$ (invariants of the regular representation) is
$1$-dimensional over $K$, spanned by $\sum_{g \in W} g$.

Therefore
$\dim_K\big(L \otimes_K (L\otimes V)^W\big) = [L:K]\cdot\dim_{\mathbb{C}}V
= |W|\dim_{\mathbb{C}}V$ by (F2), which equals
$\dim_K(L\otimes_{\mathbb{C}}V) = [L:K]\dim_{\mathbb{C}}V$ as well.
An injective $K$-linear map between $K$-vector spaces of the same
finite dimension is an isomorphism; since $\mu_V$ is visibly
$L$-linear and injective as a $K$-linear map, it is bijective, hence
an isomorphism of $L$-vector spaces.
\end{proof}

\begin{lemma}[The extended partials span $\operatorname{Der}_{\mathbb{C}}(L)$]\label{lem:der-span}
The $L$-vector space $\operatorname{Der}_{\mathbb{C}}(L)$ of
$\mathbb{C}$-linear derivations of $L$ is free of rank $n$, with basis
given by the unique extensions
$\widetilde{\partial/\partial y_1},\dots,\widetilde{\partial/\partial y_n}$
of the derivations $\partial/\partial y_1,\dots,\partial/\partial y_n$
of $K$ furnished by (F1).
\end{lemma}

\begin{proof}
Since $K = \mathbb{C}(y_1,\dots,y_n)$ is purely transcendental,
$\operatorname{Der}_{\mathbb{C}}(K)$ is a free $K$-vector space of
rank $n$ with basis $\partial/\partial y_1,\dots,\partial/\partial y_n$;
equivalently, the module of K\"ahler differentials
$\Omega_{K/\mathbb{C}}$ (see e.g.\ \cite[\S 25]{Matsumura1989} for
the general theory used throughout this proof)
is free of rank $n$ with dual basis $dy_1,\dots,dy_n$.

Since $L/K$ is a finite extension of characteristic-$0$ fields, it is
separable, hence \emph{unramified}: $\Omega_{L/K} = 0$. The conormal
exact sequence for $\mathbb{C} \to K \to L$,
\[
  \Omega_{K/\mathbb{C}} \otimes_K L \longrightarrow \Omega_{L/\mathbb{C}}
    \longrightarrow \Omega_{L/K} = 0,
\]
therefore shows that $\Omega_{K/\mathbb{C}} \otimes_K L \to \Omega_{L/\mathbb{C}}$
is surjective. Both sides are finite-dimensional $L$-vector spaces:
the left side has dimension $n$ (base change of a rank-$n$ free
module), and the right side has dimension
$\operatorname{trdeg}_{\mathbb{C}}(L) = \operatorname{trdeg}_{\mathbb{C}}(K) = n$
(transcendence degree is unchanged under the finite extension $L/K$,
and equals $\dim_L \Omega_{L/\mathbb{C}}$ since $L/\mathbb{C}$ is
finitely generated and separably generated in characteristic $0$). A
surjective linear map between $L$-vector spaces of the same finite
dimension $n$ is an isomorphism, so
$\Omega_{K/\mathbb{C}} \otimes_K L \xrightarrow{\ \sim\ } \Omega_{L/\mathbb{C}}$.
Dualizing ($\operatorname{Der}_{\mathbb{C}} = \operatorname{Hom}_{(-)}(\Omega_{(-)/\mathbb{C}}, -)$,
see e.g.\ \cite[\S 25--26]{Matsumura1989})
gives
$\operatorname{Der}_{\mathbb{C}}(L) \cong \operatorname{Der}_{\mathbb{C}}(K) \otimes_K L$.
Explicitly, this isomorphism is induced by the canonical base-change
map $dy_j \otimes 1 \mapsto dy_j$ on differentials (the same symbol
$dy_j$, computed once in $K$ and once in $L$); dualizing, it sends
$\partial/\partial y_j \otimes 1$ to the derivation of $L$ agreeing
with $\partial/\partial y_j$ on each generator $y_k$ of $K$ (i.e.\
$dy_k \mapsto \delta_{jk}$) -- which is precisely the characterizing
property of the unique extension $\widetilde{\partial/\partial y_j}$
of (F1), two derivations of $L$ that agree on the generators $y_k$ of
$K$ and both extend a $\mathbb{C}$-derivation necessarily coinciding,
by that same uniqueness. Hence the isomorphism is exactly the
statement that
$\widetilde{\partial/\partial y_1},\dots,\widetilde{\partial/\partial y_n}$
form an $L$-basis of $\operatorname{Der}_{\mathbb{C}}(L)$.
\end{proof}

\begin{lemma}[Rigidity of flat endomorphisms]\label{lem:rigidity}
Let $V$ be a finite-dimensional $\mathbb{C}$-vector space and equip
$L \otimes_{\mathbb{C}} V$ with the connection
$\widetilde{D} \otimes \operatorname{id}$ for $\widetilde{D}$ ranging
over derivations of $L$ extending those of $K$. Then any $L$-linear
endomorphism of $L \otimes_{\mathbb{C}} V$ commuting with
$\widetilde{D} \otimes \operatorname{id}$ for all such $\widetilde{D}$
lies in $\operatorname{id}_L \otimes \operatorname{End}_{\mathbb{C}}(V)$.
\end{lemma}

\begin{proof}
Fix a $\mathbb{C}$-basis $v_1,\dots,v_m$ of $V$ and write an
$L$-linear endomorphism as a matrix $(f_{ij}) \in \operatorname{Mat}_m(L)$
acting by $v_j \mapsto \sum_i f_{ij} v_i$ (extended $L$-linearly), so
that $\Phi(a \otimes v_j) = \sum_i a f_{ij} \otimes v_i$ for $a \in L$.
Fix a derivation $\widetilde D$ of $L$ extending one of $K$. Computing
$\Phi \circ (\widetilde D \otimes \operatorname{id})$ and
$(\widetilde D \otimes \operatorname{id}) \circ \Phi$ on $a \otimes v_j$
and comparing, using the Leibniz rule
$\widetilde D(af_{ij}) = \widetilde D(a) f_{ij} + a\, \widetilde D(f_{ij})$:
\[
  (\widetilde D \otimes \operatorname{id})\big(\Phi(a \otimes v_j)\big)
  = \sum_i \big[\widetilde D(a) f_{ij} + a\, \widetilde D(f_{ij})\big] \otimes v_i,
  \qquad
  \Phi\big((\widetilde D \otimes \operatorname{id})(a \otimes v_j)\big)
  = \sum_i \widetilde D(a)\, f_{ij} \otimes v_i.
\]
Commutativity of $\Phi$ with $\widetilde D \otimes \operatorname{id}$
means these two expressions agree for every $a \in L$; subtracting
leaves $\sum_i a\, \widetilde D(f_{ij}) \otimes v_i = 0$ for every
$a \in L$, and since the $v_i$ are linearly independent, taking
$a = 1$ already forces $\widetilde D(f_{ij}) = 0$ for every $i,j$.
Commuting with $\widetilde D \otimes \operatorname{id}$ for every such
$\widetilde D$ therefore forces
$\widetilde D(f_{ij}) = 0$ for every such $\widetilde D$ and every
$i,j$; by Lemma~\ref{lem:der-span}, these $\widetilde D$ already span
all of $\operatorname{Der}_{\mathbb{C}}(L)$ over $L$, so in fact
$D(f_{ij}) = 0$ for \emph{every} $\mathbb{C}$-derivation $D$ of $L$.
An element of $L$ annihilated by every derivation of $L$ over
$\mathbb{C}$ must be algebraic over $\mathbb{C}$ -- otherwise it would
be part of a transcendence basis and some derivation could be chosen
to act nontrivially on it -- hence lies in $\mathbb{C}$ itself, since
$\mathbb{C}$ is algebraically closed. Thus $f_{ij} \in \mathbb{C}$ for
all $i,j$, i.e.\ the endomorphism lies in
$\mathbb{C} \otimes \operatorname{End}(V) = \operatorname{End}_{\mathbb{C}}(V)$.
\end{proof}

\begin{remark}
A proof of Theorem~\ref{thm:descent} can be given via exactly the
three ingredients recorded above: additivity/exactness of $\nabla$
(from flat base change and Maschke's theorem), the descent isomorphism
of Lemma~\ref{lem:descent-iso} (to reduce full faithfulness to a
statement about $L \otimes_{\mathbb{C}} V$), and the rigidity of
Lemma~\ref{lem:rigidity} (to pin down $\mathbb{C}$-linear maps
commuting with the connection); this is essentially how
\cite[\S 2.4]{Nonkane2019} establishes the general
statement, though we have not reproduced the argument there and
cannot vouch for an exact correspondence. We do not spell out the
argument here; we record these ingredients because
Corollaries~\ref{cor:twisted-invariants}
and~\ref{cor:characterization} below use them directly, independently
of Theorem~\ref{thm:descent} itself.
\end{remark}

\begin{corollary}[Twisted-invariants description of the simple summands]\label{cor:twisted-invariants}
For every $\lambda \in \mathcal{P}_{r,n}/\sim$ and $0 \le l < e(\lambda)$,
\[
  N_\lambda := \widetilde{\mathcal{D}}_Y F^l_{\lambda,1}
  \;\cong\; \nabla\big(V^l(\lambda)\big)
  \;=\; \big(\widetilde{\mathcal{O}}_X \otimes_{\mathbb{C}} V^l(\lambda)\big)^{W}
\]
as $\widetilde{\mathcal{D}}_Y$-modules.
\end{corollary}

\begin{proof}
We first show $\nabla(V^l(\lambda))$ is simple. Let $M \subset \nabla(V^l(\lambda))$
be a nonzero $\widetilde{\mathcal{D}}_Y$-submodule and set
$\overline M := L \otimes_{\widetilde{\mathcal{O}}_Y} M$, a nonzero
$L$-subspace of
$L \otimes_{\widetilde{\mathcal{O}}_Y} \nabla(V^l(\lambda)) \cong L \otimes_{\mathbb{C}} V^l(\lambda)$
(Lemma~\ref{lem:descent-iso}) that is stable under
$\widetilde D \otimes \operatorname{id}$ for every derivation
$\widetilde D$ of $L$ extending one of $K$, since $M$ is a
$\widetilde{\mathcal{D}}_Y$-submodule. Write $V := V^l(\lambda)$ and
fix a $\mathbb{C}$-basis $v_1,\dots,v_f$ of $V$. Pick
$0 \ne w = \sum_i a_i \otimes v_i \in \overline M$ with a minimal
number $r$ of nonzero coefficients among all nonzero elements of
$\overline M$; relabel so $a_1,\dots,a_r \ne 0$, and scale (using that
$\overline M$ is an $L$-subspace) so $a_1 = 1$. For any extended
derivation $\widetilde D$, stability gives
$\widetilde D(w) = \sum_i \widetilde D(a_i) \otimes v_i \in \overline M$;
since $\widetilde D(a_1) = \widetilde D(1) = 0$, this element has at
most $r-1$ nonzero coefficients, so by minimality of $r$ it must be
$0$ -- exactly the minimal-support argument of
Lemma~\ref{lem:descent-iso}'s injectivity proof, now applied to a
subspace rather than a relation. Hence $\widetilde D(a_i) = 0$ for
every $i$ and every such $\widetilde D$; by Lemma~\ref{lem:der-span}
these span $\operatorname{Der}_{\mathbb{C}}(L)$ over $L$, so
$D(a_i) = 0$ for \emph{every} $\mathbb{C}$-derivation $D$ of $L$, and
as in the proof of Lemma~\ref{lem:rigidity} this forces $a_i \in \mathbb{C}$
for every $i$. Thus $w = 1 \otimes u$ for the nonzero vector
$u := \sum_i a_i v_i \in V$.

Let $U := \{u \in V : 1 \otimes u \in \overline M\}$, a
$\mathbb{C}$-subspace with $U \ne 0$ by the above; clearly
$L \otimes_{\mathbb{C}} U \subset \overline M$. If this inclusion were
strict, the same argument applied to a nonzero element of
$\overline M$ chosen outside $L \otimes_{\mathbb{C}} U$ (with minimal
support relative to a basis of $V$ extending one of $U$) would
produce a further nonzero constant vector in $U$ not already
accounted for, contradicting finite-dimensionality after finitely
many repetitions; hence $\overline M = L \otimes_{\mathbb{C}} U$.
Since $V$ is a simple $\mathbb{C}[W]$-module and $M$, hence
$\overline M$, hence $U$, is $W$-stable (as $M$ is a
$\widetilde{\mathcal{D}}_Y$-submodule of $\nabla(V)$, and the
identification with $L\otimes_{\mathbb C}V$ is $W$-equivariant), $U$
is a nonzero $W$-submodule of $V$, so $U = V$. Thus
$\overline M = L \otimes_{\mathbb{C}} V$, i.e.\ $M$ has full rank
$\dim_{\mathbb{C}} V$ inside $\nabla(V)$, which itself has rank
$\dim_{\mathbb{C}} V$ (Lemma~\ref{lem:descent-iso}); a submodule of
full rank in a torsion-free module of the same rank, with both sides
finitely generated over the same ring, forces $M = \nabla(V)$.
Hence $\nabla(V^l(\lambda))$ is simple.

The $\widetilde{\mathcal{O}}_Y$-rank of
$\nabla(V^l(\lambda))$ equals $\dim_{\mathbb{C}} V^l(\lambda) = f^l_\lambda$,
matching the rank of $N_\lambda$ (Theorem~\ref{thm:main-general}(iii),
Proposition~\ref{prop:decomposition-general}), and both are simple
$\widetilde{\mathcal{D}}_Y$-modules whose associated $W$-representation
(via $\operatorname{Loc}$) is $V^l(\lambda)$. By the uniqueness in
Theorem~\ref{thm:main-general}(ii)--(iii), $N_\lambda \cong \nabla(V^l(\lambda))$.
\end{proof}

Corollary~\ref{cor:twisted-invariants} gives an \emph{explicit,
generator-free} construction of $N_\lambda$: it is literally the
$W$-invariants of $\widetilde{\mathcal{O}}_X \otimes V^l(\lambda)$,
exactly as one would construct the vector bundle (or local system)
associated to a representation on a Galois cover -- in contrast with
Theorem~\ref{thm:main-general}(iii), which identifies $N_\lambda$ only
via an explicit but ad hoc generator, the higher Specht polynomial
$F^l_{S,T}$.

\begin{example}[The functor $\nabla$ on $D_2 = G(2,2,2)$]\label{ex:nabla-D2}
Let $L = \mathbb{C}(x_1,x_2)$, $K = \mathbb{C}(y_1,y_2)$, and
$W = D_2 = \{\mathrm{id},\sigma,\tau,\sigma\tau\}$ as in
Subsection~\ref{subsec:D2-example} ($\sigma$ the swap, $\tau$ the
simultaneous sign change). For a character
$\chi : D_2 \to \{\pm 1\}$ and $V_\chi = \mathbb{C}v_\chi$, the
diagonal action on $L \otimes V_\chi$ is
$g\cdot(a\otimes v_\chi) = g(a)\chi(g)\otimes v_\chi$, so
\[
  \nabla(\chi) = \big(L\otimes V_\chi\big)^{D_2}
  = \{a \in L \mid g(a) = \chi(g)\,a \ \ \forall g \in D_2\}
\]
is exactly the classical space of \emph{relative invariants} (or
semi-invariants) of character $\chi$ -- a $1$-dimensional $K$-vector
space, since $\dim_{\mathbb{C}} V_\chi = 1$.

The four characters of $D_2$ are realized by:
\[
  \begin{array}{c|c|c}
  \chi & (\chi(\sigma),\chi(\tau)) & \text{generator } F_\chi \\\hline
  \mathrm{triv} & (+1,+1) & 1 \\
  \mathrm{sign} & (-1,+1) & x_1^2 - x_2^2 \\
  \chi^0 & (+1,-1) & x_1+x_2 \\
  \chi^1 & (-1,-1) & x_1-x_2
  \end{array}
\]
A direct check confirms each: e.g.\ $\sigma(x_1^2-x_2^2)=x_2^2-x_1^2=-(x_1^2-x_2^2)$
and $\tau(x_1^2-x_2^2) = (-x_1)^2-(-x_2)^2 = x_1^2-x_2^2$, matching
$(\chi(\sigma),\chi(\tau))=(-1,+1)$. Since each $F_\chi$ satisfies
$g(F_\chi)=\chi(g)F_\chi$, it lies in $\nabla(\chi)$; as
$\dim_K \nabla(\chi)=1$, we get $\nabla(\chi) = K\cdot F_\chi$. These
are exactly the higher Specht polynomials $F,F',F^0,F^1$ computed by
an entirely different route in Subsection~\ref{subsec:D2-example}
(Young symmetrizers and the shift construction), illustrating
Corollary~\ref{cor:twisted-invariants} concretely.

\emph{Full faithfulness in action.} Since $\mathrm{Hom}_{D_2}(\mathrm{triv},\mathrm{sign})=0$,
Theorem~\ref{thm:descent} predicts
$\mathrm{Hom}_{\widetilde{\mathcal{D}}_Y}(\nabla(\mathrm{triv}),\nabla(\mathrm{sign}))=0$.
Concretely: a morphism $K\cdot 1 \to K\cdot(x_1^2-x_2^2)$ would be
multiplication by some $c \in L$ with $c \in K\cdot(x_1^2-x_2^2)$ and
commuting with $\partial_{y_1},\partial_{y_2}$; by
Lemma~\ref{lem:rigidity}, $c$ must be constant, but no nonzero
constant multiple of $1$ lies in $K\cdot(x_1^2-x_2^2)$, since
$x_1^2 - x_2^2 \notin K = \mathbb{C}(y_1,y_2)$. Hence
$\mathrm{Hom} = 0$, as predicted.

\emph{Where this saves work.} For the $2$-dimensional representation
$V_{(2,1)}$ of $D_3$, computing $(L\otimes V_{(2,1)})^{D_3}$ directly
(vector-valued semi-invariants, a $2\times 2$ matrix computation) is
considerably more involved than the tableau computation of
Subsection~\ref{subsec:D3-example}. Corollary~\ref{cor:twisted-invariants}
guarantees for free that
$\nabla(V_{(2,1)}) \cong N_{(2,1)} = \widetilde{\mathcal{D}}_Y \cdot G_1$
with $G_1 = x_1^2-x_3^2$ already computed there -- the categorical
equivalence transports a harder computation (vector-valued
semi-invariants) onto one already in hand (higher Specht polynomials),
at no additional cost.
\end{example}

\begin{corollary}[Categorical characterization of the summands of $\widetilde{\mathcal{O}}_X$]\label{cor:characterization}
A finitely generated $\widetilde{\mathcal{D}}_Y$-module $M$ is
isomorphic to a direct summand of $\widetilde{\mathcal{O}}_X$ if and
only if $\widetilde{\mathcal{O}}_X \otimes_{\widetilde{\mathcal{O}}_Y} M$
is trivial in the sense of Theorem~\ref{thm:descent}: isomorphic to
$\widetilde{\mathcal{O}}_X \otimes_{\mathbb{C}} U$ for some
finite-dimensional $U$, with $W$ acting only through the
$\widetilde{\mathcal{O}}_X$ factor.
\end{corollary}

\begin{proof}
($\Rightarrow$) By Proposition~\ref{prop:decomposition-general}, $M$ is a
direct sum of copies of various $N_\lambda \cong \nabla(V^l(\lambda))$;
base-changing along $\widetilde{\mathcal{O}}_X \otimes_{\widetilde{\mathcal{O}}_Y}(-)$
undoes the descent, giving $\widetilde{\mathcal{O}}_X \otimes V^l(\lambda)$
with $W$ acting only through the $\widetilde{\mathcal{O}}_X$ factor --
precisely the triviality condition just recalled, via the same
untwisting computation as Lemma~\ref{lem:descent-iso}, now applied at
the level of $\widetilde{\mathcal{O}}_X$ rather than $L$.
($\Leftarrow$) This is exactly membership in the essential image of
$\nabla$ described in Theorem~\ref{thm:descent}, and every object in
that essential image is (by Corollary~\ref{cor:twisted-invariants} and
additivity of $\nabla$) a direct sum of $N_\lambda$'s, hence a direct
summand of $\widetilde{\mathcal{O}}_X$ by
Proposition~\ref{prop:decomposition-general}.
\end{proof}

\begin{remark}
To be precise about attribution: Theorem~\ref{thm:descent} itself is
due to \cite[\S 2.4]{Nonkane2019} (for a single symmetric group),
recalled and applied to products of symmetric groups
in \cite[Prop.\ 3.10]{NonkaneTodjihounde2023}, and specialized above
to the case $W=G(r,p,n)$; Corollaries~\ref{cor:twisted-invariants}
and~\ref{cor:characterization} do not appear in
\cite{Nonkane2019,NonkaneTodjihounde2023} and are, to our knowledge,
new. They provide a route to recognizing membership in the
decomposition of Proposition~\ref{prop:decomposition-general} that
bypasses the explicit construction of a generating higher Specht
polynomial.
\end{remark}

\section{The real case: $W(D_n) = G(2,2,n)$}\label{sec:Dn}

We now specialize Section~\ref{sec:general} to $r=p=2$, recovering the
real reflection group $W(D_n)$ -- our own starting point for this
generalization, first treated by an independent, group-specific
argument in~\cite{NonkaneLawson2021}. We know that $W(D_n) = G(2,2,n)$: let
$\epsilon_1,\dots,\epsilon_n$ be the standard basis of
$\h = \mathbb{R}^n$. The corresponding set of reflections
$\{\alpha_s \mid s \in S\}$ consists of the $n(n-1)$ roots
$\epsilon_i - \epsilon_j$ and $\epsilon_i + \epsilon_j$
($1 \le i < j \le n$) -- one reflection for each sign choice per
pair, matching the general count of Section~\ref{sec:general}
(non-diagonal reflections only, since $d=1$ here).

\begin{corollary}[Discriminant, invariants, and module structure]\label{cor:Dn-setup}
For $W = W(D_n) = G(2,2,n)$ (so $r=2$, $p=2$, $d=1$), the fundamental
invariants are
\[
  y_j = \sum_{i=1}^n x_i^{2j}, \quad 1 \le j \le n-1, \qquad y_n = x_1\cdots x_n,
\]
and by Lemma~\ref{lem:jacobian-general},
\[
  \Delta = \det A = (-1)^{n-1}\, 2^{\,n-1}(n-1)! \prod_{1 \le i < j \le n} (x_j^2 - x_i^2).
\]
In particular $\det A \ne 0$ on $\widetilde{\mathcal{O}}_X$, so
$\widetilde{\mathcal{O}}_X$ is a $\widetilde{\mathcal{D}}_Y$-module.
\end{corollary}

\begin{proof}
Immediate from Lemma~\ref{lem:jacobian-general} with $r=2,p=2,d=1$
(and $n$ arbitrary): the second product in the general formula is
empty when $d=1$, the prefactor
$(-1)^{n-1}dr^{n-1}(n-1)!=(-1)^{n-1}2^{n-1}(n-1)!$, and the remaining
data match $D_n$'s classical root system
$\epsilon_i-\epsilon_j,\ \epsilon_i+\epsilon_j$. (As a check, $n=2$
gives $\Delta=-2(x_2^2-x_1^2)=2(x_1^2-x_2^2)$, matching the direct
computation of Subsection~\ref{subsec:D2-example}.)
\end{proof}

Combining Corollary~\ref{cor:Dn-setup} with Theorem~\ref{thm:main-general}
and Proposition~\ref{prop:decomposition-general} (specialized to
$r=p=2$, so $\mathcal{P}_{r,n}/\!\sim\, = \mathcal{P}_{2,n}/\!\sim$ as
in Subsection~\ref{subsec:specht}) immediately gives the decomposition theorem for
$W(D_n)$:

\begin{corollary}[Decomposition theorem for $W(D_n)$]\label{cor:Dn-main}
For every primitive idempotent $e \in \mathbb{C}[W(D_n)]$:
\begin{enumerate}[label=(\roman*)]
  \item $e\widetilde{\mathcal{O}}_X$ is a nontrivial
    $\widetilde{\mathcal{D}}_Y$-submodule of $\widetilde{\mathcal{O}}_X$;
  \item the $\widetilde{\mathcal{D}}_Y$-module $e\widetilde{\mathcal{O}}_X$
    is simple;
  \item there exist $\lambda \in \mathcal{P}_{2,n}/\sim$ and a higher
    Specht polynomial $F^l_{S,T}$ such that
    $e\widetilde{\mathcal{O}}_X = \widetilde{\mathcal{D}}_Y F^l_{S,T}$.
\end{enumerate}
Moreover,
\[
  \widetilde{\mathcal{O}}_X = \bigoplus_{\lambda \in \mathcal{P}_{2,n}/\sim}
    \bigoplus_{l=0}^{e(\lambda)-1} f^l_\lambda\,
    \widetilde{\mathcal{D}}_Y F^l_{\lambda,1},
  \qquad f^l_\lambda = \dim_{\mathbb{C}}\big(V^l(\lambda)\big).
\]
\end{corollary}

This considerably shortens what a self-contained treatment of
$W(D_n)$ alone would require: an independent six-step argument (a
Jacobian lemma, an embedding proposition, an isomorphism corollary,
the main structure theorem, a further isomorphism corollary, and a
decomposition proposition, mirroring
Lemma~\ref{lem:jacobian-general} through
Proposition~\ref{prop:decomposition-general} for the general case) is
replaced here by a two-line corollary, since none of the general
argument uses any feature of $D_n$ beyond Corollary~\ref{cor:Dn-setup}.

We now illustrate Corollary~\ref{cor:Dn-main} on the two smallest
cases, $n=2$ and $n=3$.

\subsection{Example: $D_2 = G(2,2,2)$}\label{subsec:D2-example}

We now work out explicitly the case $n=2$ of $W(D_n)$. Since $|D_2|=4$, $D_2 \cong \mathbb{Z}/2\times\mathbb{Z}/2$ is
abelian, generated by the transposition
$\sigma:(x_1,x_2)\mapsto(x_2,x_1)$ and the simultaneous sign change
$\tau:(x_1,x_2)\mapsto(-x_1,-x_2)$; all four irreducible
representations are one-dimensional.

Invariants: $y_1=x_1^2+x_2^2$, $y_2=x_1x_2$ (both of degree $2$), and
$\Delta = 2(x_1^2-x_2^2)$, matching Corollary~\ref{cor:Dn-setup} with $n=2$.

Applying the shift construction of Subsection~\ref{subsubsec:shift} to $G(2,2,2)$
yields the following basis of $\Lambda_0 \cong \mathbb{C}[D_2]$, with
standard tableaux as indicated (component $0$ only for $F,F'$; both
components for $F^0,F^1$, as in Subsection~\ref{subsec:Gr1n}):
\[
 F = 1 \ \big[\begin{array}{|c|c|}\hline 1&2\\\hline\end{array}\big],
 \qquad
 F' = x_1^2-x_2^2 \ \big[\begin{array}{|c|}\hline 1\\\hline 2\\\hline\end{array}\big],
\]
\[
 F^0 = x_1+x_2 \ \Big[\big(0:\begin{array}{|c|}\hline 1\\\hline\end{array},\
   1:\begin{array}{|c|}\hline 2\\\hline\end{array}\big)\Big],
 \qquad
 F^1 = x_1-x_2 \ \Big[\big(0:\begin{array}{|c|}\hline 2\\\hline\end{array},\
   1:\begin{array}{|c|}\hline 1\\\hline\end{array}\big)\Big],
\]
of respective degrees $0,2,1,1$, and respective characters
$(\sigma,\tau) \mapsto (+1,+1),\ (-1,+1),\ (+1,-1),\ (-1,-1)$ -- the
four distinct characters of $\mathbb{Z}/2\times\mathbb{Z}/2$.

\begin{proposition}
The Hilbert series of $\Lambda_0 = \mathcal{O}_X/J_+$ for $D_2$ is
\[
 \frac{(1-t^2)^2}{(1-t)^2} = (1+t)^2 = 1+2t+t^2,
\]
so $\dim_{\mathbb{C}} \Lambda_0 = 1+2+1 = 4 = |D_2|$, matching the
basis above (degree $0$: $F$; degree $1$: $F^0,F^1$; degree $2$: $F'$).
\end{proposition}

\begin{remark}
$D_2$ has index $2$ in $B_2 = G(2,1,2)$ ($|D_2|=4$, $|B_2|=8$, see
Subsection~\ref{subsec:Gr1n} below), and the sign character
$F'=x_1^2-x_2^2$ is, up to a scalar, exactly the discriminant
$\Delta$ -- the invariant one expects to carry the ``new'' sign
information lost upon passing from the ambient group $G(r,1,n)$ to
the subgroup $G(r,r,n)$.
\end{remark}

\subsection{Example: $D_3 = G(2,2,3)$}\label{subsec:D3-example}

We now illustrate the case $n=3$, where $W(D_3) \cong \mathfrak{S}_4$,
$|D_3|=24$, with five irreducible representations of dimensions
$1,1,2,3,3$ ($1+1+4+9+9=24$).

Invariants: $y_1=\sum x_i^2$ (degree~$2$), $y_2=\sum x_i^4$
(degree~$4$), $y_3=x_1x_2x_3$ (degree~$3$); the product of degrees is
$2\cdot4\cdot3=24=|D_3|$, as it must be.

The five equivalence classes $\lambda \in \mathcal{P}_{2,3}/\sim$ (each with
$e(\lambda)=1$, since $n=3$ is odd and no shape is shift-invariant)
correspond to the five irreducibles:

\begin{center}
\begin{tabular}{lcl}
shape $\lambda^1$ (untwisted component) & dim & role \\
\hline
$(3)$ & $1$ & trivial \\
$(1,1,1)$ & $1$ & sign \\
$(2,1)$ & $2$ & standard rep.\ of $\mathfrak{S}_3 \subset \mathfrak{S}_4$ \\
$(2)$ with a box in component $1$ & $3$ & standard rep.\ of $\mathfrak{S}_4$ \\
$(1,1)$ with a box in component $1$ & $3$ & standard $\otimes$ sign \\
\end{tabular}
\end{center}

\paragraph{The two one-dimensional pieces.}
Trivial: $F=1$ (degree $0$). Sign ($\lambda^1=(1,1,1)$): the full
antisymmetrizer applied to the appropriate monomial produces a
Vandermonde determinant in the $x_i^2$:
\[
 F' \;\propto\; \det\big(x_i^{2(j-1)}\big)_{1\le i,j\le3}
 = \prod_{i<j}(x_j^2-x_i^2) \;\propto\; \Delta \qquad (\deg 6).
\]
As for $D_2$ above, \emph{the sign representation of $D_n$
is always given, up to scalar, by the discriminant itself} -- this
holds for every $n$, by the same Vandermonde computation used in the
proof of Lemma~\ref{lem:jacobian-general}.

\paragraph{The dimension-2 piece, shape $(2,1)$.}
Let $T_A,T_B$ be the two standard tableaux of shape $(2,1)$ (hook
lengths $3,1,1$, so $\dim=3!/3=2$):
\[
  T_A = \begin{array}{|c|c|}\hline 1&2\\\hline 3\\\cline{1-1}\end{array},
  \qquad
  T_B = \begin{array}{|c|c|}\hline 1&3\\\hline 2\\\cline{1-1}\end{array}.
\]
Reading indices from $S=T_A$ gives
$i(S)=(0,0,1)$ on the cells (row-major), so the base monomial is
$x_3^{2}$ (only the cell of index $1$ contributes). Applying the Young
symmetrizers:
\[
 G_1 = e_{T_A}(x_3^2) \;\propto\; x_1^2 - x_3^2,
 \qquad
 G_2 = e_{T_B}(x_2^2) \;\propto\; x_1^2 - x_2^2.
\]
These are exactly the classical Specht polynomials of the standard
representation of $\mathfrak{S}_3$, rewritten in the variables
$x_i^2$ -- consistent with the fact that any representation carried
entirely by the untwisted component ($\lambda^2=\emptyset$) is
insensitive to the sign changes in $D_n$ and reduces to ordinary
Specht theory for $\mathfrak{S}_n$ in the variables $x_i^2$.

\paragraph{The two three-dimensional pieces.} These are the shapes
mixing both components, where the genuinely $D_n$-specific structure
appears; rather than revisiting the tableau machinery, we obtain both
directly.

\emph{Standard representation.} No invariant has degree $1$ (the
$y_j$ have degrees $2,3,4$), so the degree-$1$ piece of $\Lambda_0$
coincides with the degree-$1$ piece of $\mathcal{O}_X$ itself, namely
$\mathrm{span}(x_1,x_2,x_3)$: this is exactly the reflection
representation of $D_3$ on its defining coordinates, necessarily
irreducible of dimension $n=3$. Hence
\[
  H_1 := x_1 \qquad (\deg 1)
\]
generates the corresponding simple $\widetilde{\mathcal{D}}_Y$-module
(its character is $\chi_{\mathrm{std}}(\mathrm{id},T,Q,C,F) = (3,1,-1,0,-1)$,
computed directly as the trace of each generator on
$\mathrm{span}(x_1,x_2,x_3)$).

\emph{Standard $\otimes$ sign.} Since $g\cdot\Delta = \mathrm{sgn}(g)\,\Delta$
for every $g \in D_3$, the chain rule shows the gradient
$(\partial\Delta/\partial x_1,\partial\Delta/\partial x_2,\partial\Delta/\partial x_3)$
transforms via the reflection representation twisted by the sign
character -- a classical fact for Coxeter groups. Explicitly, with
$\Delta = (x_2^2-x_1^2)(x_3^2-x_1^2)(x_3^2-x_2^2)$,
\[
  H_1' := \frac{\partial \Delta}{\partial x_1}
    = -2x_1(x_2^2-x_3^2)(2x_1^2-x_2^2-x_3^2) \qquad (\deg 5).
\]
We verify this directly (rather than relying on the general fact
alone): writing $H_2' := \partial\Delta/\partial x_2$, a direct
computation gives $H_1'(x_2,x_1,x_3) = -H_2'(x_1,x_2,x_3)$ under the
swap $T$, and $H_1'(-x_1,-x_2,x_3) = -H_1'(x_1,x_2,x_3)$ under the
sign change $Q$ -- both matching
$\chi_{\mathrm{std}\otimes\mathrm{sgn}} = \chi_{\mathrm{std}}\cdot\chi_{\mathrm{sgn}}
= (3,-1,-1,0,1)$, and confirming $H_1'$ generates the remaining
simple $\widetilde{\mathcal{D}}_Y$-module.

\subsection{Central primitive idempotents for $D_2$ and $D_3$}

For a finite group $W$ and an irreducible character $\chi_\lambda$ of
dimension $\dim\lambda$, the central primitive idempotent projecting
onto the $\lambda$-isotypic component is
\[
  r_\lambda = \frac{\dim\lambda}{|W|} \sum_{g \in W} \overline{\chi_\lambda(g)}\, g.
\]
We compute these explicitly for the two examples above.

\paragraph{$D_2 = \{\mathrm{id},\sigma,\tau,\sigma\tau\}$.}
Being abelian of order $4$, all characters are $1$-dimensional, and
the idempotents are read off directly from the character table:
\[
  r_{\mathrm{triv}} = \tfrac14(\mathrm{id}+\sigma+\tau+\sigma\tau)
  \quad [F=1],
  \qquad
  r_{\mathrm{sign}} = \tfrac14(\mathrm{id}-\sigma+\tau-\sigma\tau)
  \quad [F'=x_1^2-x_2^2],
\]
\[
  r^{0} = \tfrac14(\mathrm{id}+\sigma-\tau-\sigma\tau)
  \quad [F^0=x_1+x_2],
  \qquad
  r^{1} = \tfrac14(\mathrm{id}-\sigma-\tau+\sigma\tau)
  \quad [F^1=x_1-x_2].
\]
One checks $r_{\mathrm{triv}}+r_{\mathrm{sign}}+r^0+r^1=\mathrm{id}$,
since the coefficients of $\sigma,\tau,\sigma\tau$ cancel pairwise.


\paragraph{$D_3 \cong \mathfrak{S}_4$.}
Denote by $T,Q,C,F$ (by abuse of notation, also standing for the sum
of the elements in each class) the conjugacy classes of $D_3$, with
representatives
\[
  \mathrm{id}\ (|{\cdot}|{=}1),\quad
  T: x_1\leftrightarrow x_2\ (|{\cdot}|{=}6),\quad
  Q: (x_1,x_2)\mapsto(-x_1,-x_2)\ (|{\cdot}|{=}3),
\]
\[
  C: (x_1,x_2,x_3)\mapsto(x_2,x_3,x_1)\ (|{\cdot}|{=}8),
  \qquad
  F: (x_1,x_2,x_3)\mapsto(x_2,-x_1,-x_3)\ (|{\cdot}|{=}6),
\]
corresponding respectively to the $\mathfrak{S}_4$-cycle types
$1^4,\,2\,1^2,\,2^2,\,3\,1,\,4$. Using the character table of
$\mathfrak{S}_4$:
\[
  r_{(3)} = \tfrac{1}{24}\big[\mathrm{id}+T+Q+C+F\big]
  \quad [\dim 1,\ F=1],
\]
\[
  r_{(1^3)} = \tfrac{1}{24}\big[\mathrm{id}-T+Q+C-F\big]
  \quad [\dim 1,\ F' \propto \Delta],
\]
\[
  r_{(2,1)} = \tfrac16\,\mathrm{id} + \tfrac16\, Q - \tfrac{1}{12}\, C
  \quad [\dim 2],
\]
\[
  r_{(2)+(1)} = \tfrac18\big[3\,\mathrm{id}+T-Q-F\big]
  \quad [\dim 3],
  \qquad
  r_{(1,1)+(1)} = \tfrac18\big[3\,\mathrm{id}-T-Q+F\big]
  \quad [\dim 3].
\]
One checks $r_{(3)}+r_{(1^3)}+r_{(2,1)}+r_{(2)+(1)}+r_{(1,1)+(1)}=\mathrm{id}$:
e.g.\ for the coefficient of $T$, $\tfrac1{24}-\tfrac1{24}+0+\tfrac18-\tfrac18=0$;
similarly for $Q$ and $F$, while $C$ only appears in $r_{(3)},r_{(1^3)},r_{(2,1)}$
with coefficients $\tfrac1{24}+\tfrac1{24}-\tfrac1{12}=0$.

\subsection{Illustrating the decomposition for $D_2$ and $D_3$}

We now instantiate Corollary~\ref{cor:Dn-main} (equivalently
Proposition~\ref{prop:decomposition-general} with $r=p=2$) on the two
examples above, exhibiting $\widetilde{\mathcal O}_X$ as an explicit
direct sum of simple $\widetilde{\mathcal D}_Y$-modules.

\paragraph{$D_2$.} All four irreducible representations of $D_2$ are
$1$-dimensional (so $f_\lambda=1$ throughout, and parts~(i) and~(ii)
of Proposition~\ref{prop:decomposition-general} coincide). Using the
generators of Subsection~\ref{subsec:D2-example},
\[
  \widetilde{\mathcal O}_X =
  \widetilde{\mathcal D}_Y\cdot 1 \;\oplus\;
  \widetilde{\mathcal D}_Y\cdot(x_1^2-x_2^2) \;\oplus\;
  \widetilde{\mathcal D}_Y\cdot(x_1+x_2) \;\oplus\;
  \widetilde{\mathcal D}_Y\cdot(x_1-x_2),
\]
four summands of $\widetilde{\mathcal O}_Y$-rank $1$ each, totalling
rank $4 = |D_2|$ -- matching, term by term, the four primitive central
idempotents $r_{\mathrm{triv}}, r_{\mathrm{sign}}, r^0, r^1$ computed
above.

\paragraph{$D_3$.} Here $D_3 \cong \mathfrak{S}_4$ has representations
of dimensions $1,1,2,3,3$. Using the explicit generators
$F=1$, $F'\propto\Delta$, $G_1=x_1^2-x_3^2$ (together with
$G_2=x_1^2-x_2^2$, isomorphic to $G_1$ by
Corollary~\ref{cor:iso-simple}), and the two generators $H_1=x_1$,
$H_1'=\partial\Delta/\partial x_1$ of Subsection~\ref{subsec:D3-example},
Proposition~\ref{prop:decomposition-general}(ii) gives
\[
  \widetilde{\mathcal O}_X =
  \widetilde{\mathcal D}_Y\cdot 1 \;\oplus\;
  \widetilde{\mathcal D}_Y\cdot \Delta \;\oplus\;
  2\,\widetilde{\mathcal D}_Y\cdot G_1 \;\oplus\;
  3\,\widetilde{\mathcal D}_Y\cdot H_1 \;\oplus\;
  3\,\widetilde{\mathcal D}_Y\cdot H_1',
\]
this time with \emph{every} summand generated explicitly. As a check
on ranks, each block contributes $f_\lambda \times f_\lambda = f_\lambda^2$
to the total $\widetilde{\mathcal O}_Y$-rank, and
\[
  1^2 + 1^2 + 2^2 + 3^2 + 3^2 = 1+1+4+9+9 = 24 = |D_3|.
\]

\section{Other particular cases}\label{sec:other-cases}

\subsection{The case $G(r,r,n)$}\label{subsec:Grrn}

The case $p=r$ (hence $d=1$) deserves special mention: it is the
family directly containing $W(D_n)=G(2,2,n)$. Here there are no
diagonal reflections, and
\[
 y_j = \sum_{i=1}^n x_i^{\,jr},\ 1\le j\le n-1, \qquad y_n = x_1\cdots x_n,
 \qquad
 \Delta_{G(r,r,n)} = c_0\prod_{1\le i<j\le n}(x_i^r-x_j^r).
\]
Lemma~\ref{lem:jacobian-general} simplifies to
\[
 \det A = (-1)^{n-1} r^{\,n-1}(n-1)! \prod_{i<j}(x_j^r-x_i^r).
\]

\begin{corollary}\label{cor:Grrn}
For every primitive idempotent $e \in \mathbb{C}[G(r,r,n)]$, the
$\widetilde{\mathcal{D}}_Y$-module $e\,\widetilde{\mathcal{O}}_X$ is simple, generated by a
higher Specht polynomial $F^l_{S,T}(\lambda)$ as in
Theorem~\ref{thm:main-general}, and
\[
 \widetilde{\mathcal{O}}_X = \bigoplus_{\lambda \in \mathcal{P}_{r,n}/\sim}
 \bigoplus_{l=0}^{e(\lambda)-1} f^l_\lambda \,
 \widetilde{\mathcal{D}}_Y F^l_{\lambda,1}.
\]
Taking $r=2$ recovers exactly Corollary~\ref{cor:Dn-main} of Section~\ref{sec:Dn}
for $W(D_n)$.
\end{corollary}

\subsection{The real case $W(B_n) = G(2,1,n)$}\label{subsec:Bn}

We single out $r=2$, $p=1$ (so $d=2$) as a second real reflection
group alongside $W(D_n)$: the Weyl group $W(B_n)$ of type $B_n$
(equivalently $C_n$), the group of all signed permutations of
$\{1,\dots,n\}$, of order $|W(B_n)| = 2^n n!$. Unlike $W(D_n)$, which
imposes an even number of sign changes, $W(B_n) = G(2,1,n)$ is the
\emph{full} hyperoctahedral group, with no parity constraint -- and
correspondingly $p=1$ here rather than $p=r$.

Its root system consists of the $n$ short roots $\epsilon_i$
($1 \le i \le n$, giving the reflections $x_i \mapsto -x_i$) together
with the $n(n-1)$ long roots $\epsilon_i \pm \epsilon_j$ ($i<j$, the
same reflections as for $D_n$); in particular
$|W(B_n)| = 2^n n!$ has exactly $n + n(n-1) = n^2$ reflections. By
Subsection~\ref{subsec:Gr1n} (taking $r=2$), the fundamental invariants
are
\[
  y_j = \sum_{i=1}^n x_i^{2j}, \quad 1 \le j \le n-1,
  \qquad
  y_n = (x_1 \cdots x_n)^2,
\]
and Lemma~\ref{lem:jacobian-general} (with $r=2$, $d=2$) gives the
discriminant
\[
  \Delta_{W(B_n)} = c_0 \Big(\prod_{i=1}^n x_i\Big) \prod_{1 \le i < j \le n} (x_j^2 - x_i^2)
\]
for some nonzero constant $c_0$ -- the classical type-$B_n$
discriminant, with one linear factor $x_i$ per short root and one
quadratic factor $x_j^2-x_i^2$ per pair of long roots.

\begin{corollary}\label{cor:Bn}
For every primitive idempotent $e \in \mathbb{C}[W(B_n)]$, the
$\widetilde{\mathcal{D}}_Y$-module $e\,\widetilde{\mathcal{O}}_X$ is
simple, generated by a higher Specht polynomial as in
Theorem~\ref{thm:main-general}, and
\[
  \widetilde{\mathcal{O}}_X = \bigoplus_{\lambda \in \mathcal{P}_{2,n}}
    f_\lambda \, \widetilde{\mathcal{D}}_Y F_{\lambda,1},
\]
the sum running over \emph{all} $\lambda \in \mathcal{P}_{2,n}$ (no
quotient by $\sim$, and no auxiliary index $l$, since $p=1$ forces
$e(\lambda)=1$ throughout -- unlike the $W(D_n)$ case, where $p=r=2$
forces the nontrivial shift identification of
Subsection~\ref{subsec:specht}).
\end{corollary}

\begin{proof}
Immediate specialization of Theorem~\ref{thm:main-general} and
Proposition~\ref{prop:decomposition-general} to $r=2$, $p=1$, exactly
as in Corollary~\ref{cor:Grrn} above.
\end{proof}

\begin{remark}
The case $n=2$ of Corollary~\ref{cor:Bn} is precisely
$G(2,1,2) = W(B_2)$, worked out fully explicitly (all five generators,
all central and primitive idempotents) in
Subsections~\ref{subsec:Gr1n}--\ref{subsec:Gr1n-idem} below, and
already used as a running example earlier in this article, in
Section~\ref{sec:general}
(Examples~\ref{ex:CW-decomp} and~\ref{ex:prop29}).
\end{remark}

\subsection{The case $G(r,1,n)$: an explicit example for $n=2$}\label{subsec:Gr1n}

We now take $p=1$ (hence $d=r$): this is the group $G(r,n)$ of
Subsection~\ref{subsec:specht}, i.e.\ the full monomial
group $(\mathbb{Z}/r\mathbb{Z})^n \rtimes S_n$, with fundamental
invariants $e_j(x_1^r,\dots,x_n^r)$, $1\le j\le n$, and no restriction
by $\sim$ (every class $\lambda \in \mathcal{P}_{r,n}$ is a singleton, since
$d=r$ forces $b(\lambda)=1$, $e(\lambda)=1$).

For $n=2$, $|G(r,1,2)| = 2r^2$. The higher Specht polynomials
(Subsection~\ref{subsec:specht}) can be written down completely; the
general-$n$ case of this decomposition is treated
in~\cite{NonkaneLawson2022}.

\paragraph{Irreducible components indexed by a single row/column,
component $\nu \in \{0,\dots,r-1\}$ (dimension 1 each).}
Shape $(2)$ in component $\nu$, standard tableau
$T = \begin{array}{|c|c|}\hline 1 & 2\\\hline\end{array}$:
\[
 F_\nu = (x_1 x_2)^{\nu}.
\]
Shape $(1,1)$ in component $\nu$, standard tableau
$T = \begin{array}{|c|}\hline 1\\\hline 2\\\hline\end{array}$:
\[
 F'_\nu = (x_1x_2)^{\nu}(x_1^{r} - x_2^{r}).
\]

\paragraph{Irreducible components indexed by a pair of components
$0\le \nu < \mu \le r-1$, each with a single box (dimension 2 each,
multiplicity 2 in the regular representation).} The two standard
tableaux are
\[
  T_a = \Big(\nu:\begin{array}{|c|}\hline 1\\\hline\end{array}
    \ ,\ \mu:\begin{array}{|c|}\hline 2\\\hline\end{array}\Big),
  \qquad
  T_b = \Big(\nu:\begin{array}{|c|}\hline 2\\\hline\end{array}
    \ ,\ \mu:\begin{array}{|c|}\hline 1\\\hline\end{array}\Big),
\]
(box $1$ in component $\nu$ or $\mu$ respectively), giving
\[
 \{x_1^{\nu}x_2^{\mu},\; x_1^{\mu}x_2^{\nu}\}
 \qquad\text{and}\qquad
 \{x_1^{r+\nu}x_2^{\mu},\; x_1^{\mu}x_2^{r+\nu}\}.
\]

\begin{proposition}\label{prop:Gr1n-basis}
The polynomials above form a basis of $\Lambda_0 \cong \mathbb{C}[G(r,1,2)]$
as a vector space, of total dimension
\[
 \underbrace{r}_{F_\nu} + \underbrace{r}_{F'_\nu}
 + \underbrace{4\binom{r}{2}}_{\text{pairs}}
 = 2r + 2r(r-1) = 2r^2 = |G(r,1,2)|.
\]
\end{proposition}

\begin{proof}
\emph{Step 1 -- Enumeration of $\mathcal{P}_{r,2}$.} An $r$-tuple
$\lambda=(\lambda^1,\dots,\lambda^r)$ of Young diagrams of total size
$2$ is of one of three types: (a) a single component $\nu$ carries
shape $(2)$, the rest empty ($r$ choices of $\nu$); (b) a single
component $\nu$ carries shape $(1,1)$ ($r$ choices); (c) two
\emph{distinct} components $\nu < \mu$ each carry shape $(1)$
($\binom{r}{2}$ choices). These exhaust $\mathcal{P}_{r,2}$, since the
only partitions of $2$ are $(2)$ and $(1,1)$, and $2=2$ or $2=1+1$ are
the only ways to split the total size across components.

\emph{Step 2 -- Cases (a)--(b).} The shapes $(2)$ and $(1,1)$ each
admit exactly one standard tableau, so $f_\lambda = 1$. Since $p=1$
gives $d=r$, the condition ``$1 \in T^\nu$ for some $\nu \le d=r$''
defining $\mathrm{STab}(\lambda)_d$ is automatic (there are only $r$
components), so $\mathrm{STab}(\lambda)_d = \mathrm{STab}(\lambda)$,
reduced to the unique tableau. This yields exactly \emph{one} higher
Specht polynomial $F^0_{S,T}$ per such $\lambda$ -- namely $F_\nu$,
resp.\ $F'_\nu$, as computed explicitly above from the construction of
Subsection~\ref{subsec:specht}.

\emph{Step 3 -- Case (c).} The two-isolated-boxes shape admits $2$
standard tableaux $T_a,T_b$ (according to whether $1$ occupies
component $\nu$ or $\mu$), so $f_\lambda = 2$, consistent with the
multinomial formula $\binom{2}{1,1}\cdot 1\cdot 1 = 2$. Again
$\mathrm{STab}(\lambda)_d = \mathrm{STab}(\lambda) = \{T_a,T_b\}$
without restriction. This yields $S \in \{T_a,T_b\}$,
$T \in \{T_a,T_b\}$, i.e.\ \emph{four} polynomials $F^0_{S,T}$ per such
$\lambda$ -- these are the two pairs
$\{x_1^\nu x_2^\mu, x_1^\mu x_2^\nu\}$ and
$\{x_1^{r+\nu}x_2^\mu, x_1^\mu x_2^{r+\nu}\}$ displayed above.

\emph{Step 4 -- Total count.}
\[
  \underbrace{r \cdot 1}_{(a)} + \underbrace{r \cdot 1}_{(b)}
  + \underbrace{\binom{r}{2}\cdot 4}_{(c)}
  = 2r + 2r(r-1) = 2r^2,
\]
which also confirms $\sum_\lambda f_\lambda^2 = |G(r,1,2)|$, as it
must (Theorem~\ref{thm:2.3}(ii)).

\emph{Step 5 -- Conclusion via Theorem~\ref{thm:2.3}.}
Theorem~\ref{thm:2.3}(iii), already proved in general in
Subsection~\ref{subsec:specht}, states that the family $\{F^l_{S,T}\}$ -- indexed over
\emph{all} $\lambda \in \mathcal{P}_{r,n}$,
$S \in \mathrm{STab}(\lambda)_d$, $T \in \mathrm{STab}(\lambda)_{db(\lambda)}$,
$l = 0,\dots,e(\lambda)-1$ -- forms a basis of the
$\mathbb{C}[x_1,\dots,x_n]^{G(r,p,n)}$-module
$\mathbb{C}[x_1,\dots,x_n]$, these elements being by construction
(Subsection~\ref{subsec:specht}) already viewed inside $\Lambda_0 = \mathcal{O}_X/J_+$.
Steps 1--3 show our list enumerates \emph{exactly} this family for
$(r,p,n)=(r,1,2)$ (where $e(\lambda)=1$ for every $\lambda$ since
$p=1$, so only $l=0$ occurs) -- with neither omission nor repetition.
Since $\Lambda_0 \cong \mathbb{C}[G(r,1,2)]$ as a vector space of
dimension $|G(r,1,2)| = 2r^2$ (Subsection~\ref{subsec:specht}), and our list has
exactly $2r^2$ elements, all drawn from the basis guaranteed by
Theorem~\ref{thm:2.3}(iii), it \emph{is} a basis of $\Lambda_0$.
\end{proof}

\begin{example}[$r=2$: the Coxeter group $B_2$]
Here $|G(2,1,2)|=8$. The basis is:
\[
 F_0 = 1,\quad F_1 = x_1x_2, \qquad
 F'_0 = x_1^2-x_2^2,\quad F'_1 = x_1x_2(x_1^2-x_2^2),
\]
together with the single pair $(\nu,\mu)=(0,1)$:
\[
 \{x_1,\, x_2\} \qquad \{x_1^2x_2,\, x_1x_2^2\}.
\]
Total: $2+2+4=8$ basis elements, matching $|G(2,1,2)|$.
\end{example}


\subsection{Central primitive idempotents for $G(r,1,2)$}\label{subsec:Gr1n-idem}

From the characters computed in Subsection~\ref{subsec:Gr1n}, we
obtain the central primitive idempotents of $G(r,1,2)$.

\paragraph{$G(r,1,2)$.}
Write elements as $(a,b;\sigma)$, $a,b \in \mathbb{Z}/r\mathbb{Z}$,
$\sigma \in \{e,\mathrm{swap}\}$, and $\xi = e^{2i\pi/r}$. Recall the
characters computed above:
\[
  \chi_{F_\nu}(a,b;\sigma) = \xi^{\nu(a+b)} \ (\forall \sigma),
  \qquad
  \chi_{F'_\nu}(a,b;\sigma) = \xi^{\nu(a+b)} \cdot
    \begin{cases} +1 & \sigma = e \\ -1 & \sigma = \mathrm{swap}, \end{cases}
\]
\[
  \chi_{(\nu,\mu)}(a,b;e) = \xi^{a\nu+b\mu}+\xi^{a\mu+b\nu},
  \qquad
  \chi_{(\nu,\mu)}(a,b;\mathrm{swap}) = 0.
\]
The corresponding central primitive idempotents are
\[
  r_{F_\nu} = \frac{1}{2r^2} \sum_{a,b=0}^{r-1} \xi^{-\nu(a+b)}
    \big[(a,b;e)+(a,b;\mathrm{swap})\big],
\]
\[
  r_{F'_\nu} = \frac{1}{2r^2} \sum_{a,b=0}^{r-1} \xi^{-\nu(a+b)}
    \big[(a,b;e)-(a,b;\mathrm{swap})\big],
\]
\[
  r_{(\nu,\mu)} = \frac{1}{r^2} \sum_{a,b=0}^{r-1}
    \big(\xi^{-a\nu-b\mu}+\xi^{-a\mu-b\nu}\big)(a,b;e), \qquad 0 \le \nu < \mu \le r-1.
\]

\begin{example}[$r=2$]
Here $\xi=-1$, and $r_{F_0} = \tfrac18\sum_{a,b\in\{0,1\}}\big[(a,b;e)+(a,b;\mathrm{swap})\big]$
is $\tfrac18$ times the sum of all $8$ group elements, while
$r_{F_1} = \tfrac18\sum_{a,b}(-1)^{a+b}\big[(a,b;e)+(a,b;\mathrm{swap})\big]$,
and $r_{(0,1)} = \tfrac14\sum_{a,b}\big[(-1)^{b}+(-1)^{a}\big](a,b;e)$,
matching the $B_2$ character table.
\end{example}

\begin{remark}
As a sanity check applicable to all three examples above ($D_2$,
$D_3$, and $G(r,1,2)$), summing all
computed central idempotents for a given group always returns the
identity element $\mathrm{id}$, and $r_\lambda r_\mu = 0$ for
$\lambda \ne \mu$ -- both are immediate from column orthogonality of
the character table.
\end{remark}

\section{Concluding remarks and perspectives}\label{sec:conclusion}

We have given a uniform construction of the simple
$\widetilde{\mathcal{D}}_Y$-module summands of the polynomial ring
localized at the discriminant, for the entire family of imprimitive
complex reflection groups $G(r,p,n)$, resting on two ingredients --
the Jacobian lemma of Subsection~\ref{subsec:jacobian} and the
double-centralizer argument of Subsection~\ref{subsec:simple-components} -- that
replace the ad hoc, group-specific constructions used in earlier
treatments of special cases (symmetric groups, and our own starting
point $W(D_n)$~\cite{NonkaneLawson2021}). Specializing this single argument recovers, and
considerably shortens, the decomposition theorems for the real
reflection groups $W(D_n)$ and $W(B_n)$, and a Galois descent
equivalence of categories (Theorem~\ref{thm:descent}) provides a
second, generator-free description of the same summands. We close by
indicating some directions this leaves open.

\paragraph{Exceptional complex reflection groups.} The Shephard--Todd
classification consists of the infinite family $G(r,p,n)$ together
with $34$ exceptional groups. The double-centralizer argument of
Subsection~\ref{subsec:simple-components} is entirely general and applies
verbatim to any complex reflection group; what is specific to
$G(r,p,n)$ is the Jacobian lemma (Lemma~\ref{lem:jacobian-general})
and the combinatorial substitute for higher Specht polynomials, both
tied to the monomial wreath-product structure. Obtaining analogous
discriminant formulas and explicit bases for the exceptional groups
would need to be developed case by case.

\paragraph{Connection to rational Cherednik algebras.} The ring
$\widetilde{\mathcal{D}}_Y$ studied here is closely related to the
spherical subalgebra of the rational Cherednik algebra $H_c(W)$ at
$t \ne 0$, $c=0$: at $c=0$ the Dunkl operators of $H_c(W)$ degenerate
to ordinary partial derivatives, and the spherical subalgebra
$eH_0(W)e$ recovers (an algebraic avatar of) the invariant
differential operator ring $D(V)^W$, closely related to
$\widetilde{\mathcal{D}}_Y$. Under this identification, our
decomposition theorem should correspond to the decomposition of the
polynomial representation of $H_0(W)$ into simple modules in category
$\mathcal{O}$ at $c=0$. Making this correspondence precise, and
comparing our explicit generators to the standard modules
$\Delta_0(\lambda)$, would connect the present combinatorial
construction to the representation theory of Cherednik algebras --
and might in turn suggest how our construction should deform as $c$
moves away from $0$.

\paragraph{Physical motivation revisited.} The original version of
this circle of ideas, for $W(D_n)$, was motivated by the rational
quantum Olshanetsky--Perelomov system; the present, purely algebraic
treatment does not use this structure. Reintroducing Dunkl--Opdam
operators for the whole family $G(r,p,n)$ -- and, in particular,
working out the associated integrable system for $W(B_n)$, now
treated on equal footing with $W(D_n)$ in
Section~\ref{sec:other-cases} -- is a natural next step.

\paragraph{Completing the remaining explicit examples.} We gave fully
explicit generators for every simple summand in the cases $D_2$,
$D_3$, and $G(r,1,2)$; extending this to $D_4$, or to a general
closed-form recipe for the generators attached to shapes mixing both
components of $\lambda$ -- rather than the case-by-case constructions
(the reflection representation and the gradient of the discriminant)
used in Subsection~\ref{subsec:D3-example} -- remains open.

\paragraph{A categorical refinement.} Corollary~\ref{cor:characterization}
identifies the summands of $\widetilde{\mathcal{O}}_X$ with a
subcategory of $\widetilde{\mathcal{D}}_Y$-modules trivialized by the
discriminant cover (Theorem~\ref{thm:descent}); it would be
interesting to know whether the equivalence of
Theorem~\ref{thm:descent} extends to a statement about \emph{all}
finitely generated $\widetilde{\mathcal{D}}_Y$-modules, in the spirit
of the Riemann--Hilbert correspondence for the discriminant complement.

\section*{Acknowledgments}

We thank Professor Rikard B\"ogvad for introducing us to this theory
and for instructive comments during the writing of this paper.

\begin{thebibliography}{99}

\bibitem{Artin1942} Artin, E., \emph{Galois Theory}, Notre Dame
Mathematical Lectures, Number 2, University of Notre Dame Press,
Notre Dame, IN, 1942.

\bibitem{ATY1997} Ariki, S., Terasoma, T., Yamada, H.,
\emph{Higher Specht polynomials}, Hiroshima Math.\ J.\ \textbf{27}
(1997), no.\ 1, 177--188.

\bibitem{Chevalley1955} Chevalley, C., \emph{Invariants of finite
groups generated by reflections}, Amer.\ J.\ Math.\ \textbf{77}
(1955), no.\ 4, 778--782.

\bibitem{FultonHarris1991} Fulton, W., Harris, J.,
\emph{Representation Theory: A First Course}, Graduate Texts in
Mathematics 129, Springer-Verlag, New York, 1991.

\bibitem{Lang2002} Lang, S., \emph{Algebra}, Revised third edition,
Graduate Texts in Mathematics 211, Springer-Verlag, New York, 2002.

\bibitem{Matsumura1989} Matsumura, H., \emph{Commutative Ring Theory},
Cambridge Studies in Advanced Mathematics 8, Cambridge University
Press, Cambridge, 1989.

\bibitem{MoritaYamada1998} Morita, W., Yamada, H., \emph{Higher
Specht polynomials for the complex reflection group $G(r,p,n)$},
Hokkaido Math.\ J.\ \textbf{27} (1998), no.\ 3, 505--515.

\bibitem{Nonkane2019} Nonkan\'e, I., \emph{Specht polynomials and
modules over the Weyl algebra}, Afr.\ Mat.\ \textbf{30} (2019), no.\
11, 279--290.

\bibitem{NonkaneLawson2021} Nonkan\'e, I., Lawson, L.M.,
\emph{Differential operators and reflection group of type $D_n$},
J.\ Phys.:\ Conf.\ Ser.\ \textbf{2090} (2021), 012098.

\bibitem{NonkaneLawson2022} Nonkan\'e, I., Lawson, L.M.,
\emph{Invariant differential operators and the generalized symmetric
group}, Gulf J.\ Math.\ \textbf{13} (2022), no.\ 2, 19--32.

\bibitem{NonkaneTodjihounde2023} Nonkan\'e, I., Todjihounde, L.,
\emph{Specht polynomials and modules over the Weyl algebra II},
J.\ Math.\ Sci.\ \textbf{271} (2023), 197--209.

\end{thebibliography}


\end{document}
